\documentclass{article}

\usepackage{amsmath, amsthm, amssymb, amsfonts}
\usepackage{thmtools}
\usepackage{graphicx}
\usepackage{setspace}
\usepackage{geometry}
\usepackage{float}
% \usepackage{hyperref}
\usepackage[utf8]{inputenc}
\usepackage[english]{babel}
\usepackage{framed}
\usepackage[dvipsnames]{xcolor}
\usepackage{tcolorbox}
\usepackage{bbold}

\colorlet{LightGray}{White!90!Periwinkle}
\colorlet{LightOrange}{Orange!15}
\colorlet{LightGreen}{Green!15}

\newcommand{\HRule}[1]{\rule{\linewidth}{#1}}

\declaretheoremstyle[name=Theorem,]{thmsty}
\declaretheorem[style=thmsty,numberwithin=section]{theorem}
\tcolorboxenvironment{theorem}{colback=LightGray}

\declaretheoremstyle[name=Proposition,]{prosty}
\declaretheorem[style=prosty,numberlike=theorem]{proposition}
\tcolorboxenvironment{proposition}{colback=LightOrange}

\declaretheoremstyle[name=Principle,]{prcpsty}
\declaretheorem[style=prcpsty,numberlike=theorem]{principle}
\tcolorboxenvironment{principle}{colback=LightGreen}

\setstretch{1.2}
\geometry{
    textheight=9in,
    textwidth=5.5in,
    top=1in,
    headheight=12pt,
    headsep=25pt,
    footskip=30pt
}

% ------------------------------------------------------------------------------

\begin{document}

% ------------------------------------------------------------------------------
% Cover Page and ToC
% ------------------------------------------------------------------------------

\title{ \normalsize \textsc{}
		\\ [2.0cm]
		\HRule{1.5pt} \\
		\LARGE \textbf{\uppercase{Full Molecular Hamiltonian}
		\HRule{2.0pt} \\ [0.6cm] \LARGE{From Classical to Quantum} \vspace*{10\baselineskip}}
		}
\date{}
\author{\textbf{Chih-En Shen}}

\maketitle
\newpage

\tableofcontents
\newpage

% ------------------------------------------------------------------------------

\section{Classical Cartesian Hamiltonian}

We have $N_{\mathrm{nucl}}$ nuclei labeled by $j$ with masses $M_j$, charges $Z_j e$, positions $\mathbf{R}_j$, and conjugate momenta $\mathbf{P}_j$; and $N_{\mathrm{elec}}$ electrons labeled by $i$ with mass $m_e$, charge $-e$, positions $\mathbf{r}_i$, and momenta $\mathbf{p}_i$. All vectors live in a space-fixed (SF) lab frame. Using Gaussian units and writing $R_{ij}=|\mathbf{R}_i-\mathbf{R}_j|$ etc.,

\begin{equation}
H = \underbrace{\sum_{j=1}^{N_{\mathrm{nucl}}}\frac{\mathbf{P}_j^{\,2}}{2M_j}}_{T_N} + \underbrace{\sum_{i=1}^{N_{\mathrm{elec}}}\frac{\mathbf{p}_i^{\,2}}{2m_e}}_{T_e} + V, \tag{1.1}
\end{equation}

\begin{equation}
V = \frac{1}{4\pi\epsilon_0}\left(\frac12\sum_{i\ne j}\frac{Z_i Z_j e^2}{R_{ij}} - \sum_{i,j}\frac{Z_j e^2}{|\mathbf{R}_j-\mathbf{r}_i|} + \frac12\sum_{i\ne j}\frac{e^2}{r_{ij}}\right). \tag{1.2}
\end{equation}

Canonical Poisson brackets $\{R_{j\alpha},P_{k\beta}\}=\delta_{jk}\delta_{\alpha\beta}$, $\{r_{i\alpha},p_{j\beta}\}=\delta_{ij}\delta_{\alpha\beta}$.

\section{Transformation to the Nuclear Center-of-Mass}

\subsection{New coordinates}

Define the nuclear COM, $3(N-1)$ internal nuclear vectors, and electron positions in the nuclear COM frame:

\begin{equation}
\mathbf{R}_{\mathrm{cm}} = \frac{1}{M_N}\sum_j M_j\mathbf{R}_j,\qquad M_N = \sum_j M_j, \tag{2.1}
\end{equation}

\begin{equation}
\mathbf{R}'_j = \mathbf{R}_j - \mathbf{R}_{\mathrm{cm}},\qquad \mathbf{r}'_i = \mathbf{r}_i - \mathbf{R}_{\mathrm{cm}}. \tag{2.2}
\end{equation}
The $\mathbf{R}'_j$ are linearly dependent:
\begin{equation}
\sum_j M_j\mathbf{R}'_j = 0. \tag{2.3}
\end{equation}

\subsection{Canonical momenta}

This is a point transformation $\mathbf{Q}_{\mathrm{new}}' = \mathbf{Q}_{\mathrm{new}}'(\mathbf{q}_{\mathrm{old}})$ with $\mathbf{q}_{\mathrm{old}} = (\mathbf{R}_j,\mathbf{r}_i)$ and $\mathbf{Q}'_{\mathrm{new}} = (\mathbf{R}_{\mathrm{cm}},\mathbf{R}'_j,\mathbf{r}'_i)$. Use the generating function:
$F_2 = F_2(\mathbf{q_{\mathrm{old}}}, \mathbf{P}, t) = \mathbf{P} \cdot \mathbf{Q}_{\mathrm{new}}'(\mathbf{q_{\mathrm{old}}})$
\begin{equation}
F_2 = \mathbf{P}_{\mathrm{cm}}\cdot\mathbf{R}_{\mathrm{cm}}(\{\mathbf{R}_j\}) + \sum_j\mathbf{P}'_j\cdot\mathbf{R}'_j(\{\mathbf{R}_j\}) + \sum_i\mathbf{p}'_{\,i}\cdot\mathbf{r}'_i(\mathbf{R}_j,\mathbf{r}_i). \tag{2.4}
\end{equation}
Then $\mathbf{P}_j = \partial F_2/\partial\mathbf{R}_j$, $\mathbf{p}_i = \partial F_2/\partial\mathbf{r}_i$. Using

\begin{equation}
\frac{\partial\mathbf{R}_{\mathrm{cm}}}{\partial\mathbf{R}_j} = \frac{M_j}{M_N}\mathbb{1},\quad \frac{\partial\mathbf{R}'_k}{\partial\mathbf{R}_j} = \bigl(\delta_{jk}-\tfrac{M_j}{M_N}\bigr)\mathbb{1},\quad \frac{\partial\mathbf{r}'_i}{\partial\mathbf{R}_j} = -\tfrac{M_j}{M_N}\mathbb{1},\quad \frac{\partial\mathbf{r}'_i}{\partial\mathbf{r}_j} = \delta_{ij}\mathbb{1}, \tag{2.5}
\end{equation}
one obtains
\begin{equation}
\mathbf{p}_i = \mathbf{p}'_{i},\qquad \mathbf{P}_j = \frac{M_j}{M_N}\bigl(\mathbf{P}_{\mathrm{cm}}-\textstyle\sum_i\mathbf{p}'_{\,i}-\sum_k\mathbf{P}'_k\bigr)+\mathbf{P}'_j. \tag{2.6}
\end{equation}
The constraint dual to (2.3) is $\sum_k\mathbf{P}'_k = 0$, so
\begin{equation}
\boxed{\;\mathbf{P}_j = \frac{M_j}{M_N}\bigl(\mathbf{P}_{\mathrm{cm}}-\textstyle\sum_i\mathbf{p}_i\bigr)+\mathbf{P}'_j,\qquad \mathbf{p}_i = \mathbf{p}'_i} \tag{2.7}
\end{equation}

Total momentum: $\sum_j\mathbf{P}_j+\sum_i\mathbf{p}_i = \mathbf{P}_{\mathrm{cm}}$. So $\mathbf{P}_{\mathrm{cm}}$ is the \emph{total} (nuclei+electrons) momentum.

\subsection{Hamiltonian after the transformation}

Insert (2.7) into $T_N$ in (1.1):

\begin{equation}
T_N = \sum_j\frac{\mathbf{P}_j^{\,2}}{2M_j} = \sum_j\frac{1}{2M_j}\left[\frac{M_j}{M_N}\bigl(\mathbf{P}_{\mathrm{cm}}-\textstyle\sum_i\mathbf{p}_i\bigr)+\mathbf{P}'_j\right]^{2}. \tag{2.8}
\end{equation}
Expanding the square and summing,
\begin{equation}
T_N = \frac{(\mathbf{P}_{\mathrm{cm}} - \sum_i\mathbf{p}_i)^2}{2M_N} + \frac{1}{M_N}\bigl(\mathbf{P}_{\mathrm{cm}}-\sum_i\mathbf{p}_i\bigr)\cdot\underbrace{\sum_j\mathbf{P}'_j}_{=0} + \sum_j\frac{\mathbf{P}_j^{'2}}{2M_j}. \tag{2.9}
\end{equation}
Combining with $T_e$ and $V$:
\begin{equation}
H = \frac{\mathbf{P}_{\mathrm{cm}}^{\,2}}{2M_N} \underbrace{- \frac{\mathbf{P}_{\mathrm{cm}}\cdot\sum_i\mathbf{p}_i}{M_N}}_{\text{artifact}} + \underbrace{\frac{(\sum_i\mathbf{p}_i)^2}{2M_N}}_{\text{mass-polarization}} + \sum_j\frac{\mathbf{P}_j^{'2}}{2M_j} + \sum_i\frac{\mathbf{p}_i^{\,2}}{2m_e} + V(\mathbf{R}'_j,\mathbf{r}'_i). \tag{2.10}
\end{equation}
Because $H$ is independent of $\mathbf{R}_{\mathrm{cm}}$, $\mathbf{P}_{\mathrm{cm}}$ is conserved; pick the eigenvalue $\mathbf{P}_{\mathrm{cm}} = 0$ (center-of-mass at rest) and drop the trivial $\mathbf{P}_{\mathrm{cm}}$-pieces:
\begin{equation}
\boxed{H_{\mathrm{int}} = \sum_j\frac{\mathbf{P}_j^{'2}}{2M_j} + \sum_i\frac{\mathbf{p}_i^{\,2}}{2m_e} + \underbrace{\frac{1}{2M_N}\Bigl(\sum_i\mathbf{p}_i\Bigr)^{2}}_{\text{electron mass polarization}} + V(\mathbf{R}'_j,\mathbf{r}'_i).}\tag{2.11}
\end{equation}
The potential is unchanged in form because $\mathbf{R}'_j-\mathbf{R}'_j = \mathbf{R}_j-\mathbf{R}_j$, $\mathbf{R}'_j-\mathbf{r}'_i = \mathbf{R}_j-\mathbf{r}_i$, $\mathbf{r}'_i-\mathbf{r}'_j = \mathbf{r}_i-\mathbf{r}_j$. The mass-polarization term, $(M_N)^{-1}\sum_{i<j}\mathbf{p}_i\cdot\mathbf{p}_j + (2M_N)^{-1}\sum_i\mathbf{p}_i^{\,2}$, is the price of referring electrons to the nuclear (rather than the total) COM. This rest-frame choice is not essential: Section 2.4 shows that the \emph{same} internal Hamiltonian follows for any $\mathbf{P}_{\mathrm{cm}}$ from a classical canonical transformation, and it is that frame-independent form we carry forward.

\subsection{Eliminating the Cross-Coupling via a Classical Canonical Transformation}

The rest-frame choice $\mathbf{P}_{\mathrm{cm}}=0$ made in Section 2.3 is convenient but stronger than necessary, and we now lift it. For $\mathbf{P}_{\mathrm{cm}}\neq 0$ the Hamiltonian (2.10) carries two artifacts inherited from having referred electrons to the \emph{nuclear} (rather than the \emph{total}) center of mass:
\begin{equation}
H = \underbrace{\frac{\mathbf{P}_{\mathrm{cm}}^{\,2}}{2M_N}}_{(A)\text{ wrong total mass}} \;\underbrace{- \frac{\mathbf{P}_{\mathrm{cm}}\!\cdot\!\sum_i\mathbf{p}_i'}{M_N}}_{(B)\text{ cross-coupling}} \;+ \frac{(\sum_i\mathbf{p}_i')^{2}}{2M_N} + \sum_j\frac{\mathbf{P}_j^{'\,2}}{2M_j} + \sum_i\frac{\mathbf{p}_i^{'\,2}}{2m_e} + V. \tag{2.12}
\end{equation}
The wrong-mass term (A) has denominator $M_N$ instead of the proper $M_{\mathrm{tot}}=M_N+N_{\mathrm{elec}}m_e$, and (B) couples the COM motion to the internal electronic momenta. We eliminate both \emph{exactly, for any $\mathbf{P}_{\mathrm{cm}}$}, by a single classical canonical transformation generated on phase space --- so that the internal Hamiltonian carried into the body-fixed treatment of Section 3 is the frame-independent object, not a rest-frame special case. This transformation is the classical preimage of the Galilean-boost unitary that one would apply in the quantum theory.

\subsubsection{Generator}

Define on phase space the function
\begin{equation}
G \;\equiv\; \frac{m_e}{M_{\mathrm{tot}}}\,\mathbf{P}_{\mathrm{cm}}\!\cdot\!\sum_i\mathbf{r}_i', \qquad M_{\mathrm{tot}} = M_N + N_{\mathrm{elec}} m_e. \tag{2.13}
\end{equation}
$G$ depends only on $\mathbf{P}_{\mathrm{cm}}$ and $\{\mathbf{r}_i'\}$; in particular it is independent of $\mathbf{R}_{\mathrm{cm}}$, $\mathbf{R}'_j$, $\mathbf{P}'_j$, and $\mathbf{p}'_i$.

\subsubsection{Induced canonical transformation}

The infinitesimal canonical transformation generated by $G$ acts on any phase-space function $f$ by $\delta f = \{f,G\}$. Because $G$ is linear in $\mathbf{P}_{\mathrm{cm}}$ and $\{\mathbf{r}'_i\}$, the bracket $\{\cdot,G\}$ applied twice annihilates the canonical variables, so the finite transformation at unit parameter is given by a single application:
\begin{equation}
\tilde f = f + \{f,G\}. \tag{2.14}
\end{equation}
Using $\{f,g\} = \sum (\partial_q f\,\partial_p g - \partial_p f\,\partial_q g)$, compute the action on each canonical pair.

\noindent\emph{Electron pair} $(\mathbf{r}'_i,\mathbf{p}'_i)$:
\begin{align*}
\{r'_{i,a},G\} \;&=\; \tfrac{\partial G}{\partial p'_{i,a}} \;=\; 0, \\
\{p'_{i,a},G\} \;&=\; -\tfrac{\partial G}{\partial r'_{i,a}} \;=\; -\tfrac{m_e}{M_{\mathrm{tot}}}P_{\mathrm{cm},a}.
\end{align*}

\noindent\emph{Nuclear-COM pair} $(\mathbf{R}_{\mathrm{cm}},\mathbf{P}_{\mathrm{cm}})$:
\begin{align*}
\{R_{\mathrm{cm},a},G\} \;&=\; \tfrac{\partial G}{\partial P_{\mathrm{cm},a}} \;=\; \tfrac{m_e}{M_{\mathrm{tot}}}\sum_i r'_{i,a}, \\
\{P_{\mathrm{cm},a},G\} \;&=\; -\tfrac{\partial G}{\partial R_{\mathrm{cm},a}} \;=\; 0.
\end{align*}
\noindent\emph{Internal nuclear pair} $(\mathbf{R}'_j,\mathbf{P}'_j)$: $G$ depends on neither, so $\{\mathbf{R}'_j,G\}=\{\mathbf{P}'_j,G\}=0$.
Hence the finite canonical transformation reads
\begin{equation}
\boxed{\;
\begin{aligned}
\tilde{\mathbf{r}}'_i &= \mathbf{r}'_i, & \tilde{\mathbf{p}}'_i &= \mathbf{p}'_i - \tfrac{m_e}{M_{\mathrm{tot}}}\mathbf{P}_{\mathrm{cm}}, \\
\tilde{\mathbf{R}}_{\mathrm{cm}} &= \mathbf{R}_{\mathrm{cm}} + \tfrac{m_e}{M_{\mathrm{tot}}}\textstyle\sum_i \mathbf{r}'_i, \quad & \tilde{\mathbf{P}}_{\mathrm{cm}} &= \mathbf{P}_{\mathrm{cm}}, \\
\tilde{\mathbf{R}}'_j &= \mathbf{R}'_j, & \tilde{\mathbf{P}}'_j &= \mathbf{P}'_j.
\end{aligned}\;}\tag{2.15}
\end{equation}

\subsubsection{Canonicity check}

The only non-obvious Poisson bracket is between the two simultaneously shifted variables:
\begin{align}
\{\tilde p'_{i,a},\tilde R_{\mathrm{cm},b}\}
&= \bigl\{\mathbf{p}'_{i,a} - \tfrac{m_e}{M_{\mathrm{tot}}}P_{\mathrm{cm},a},\; R_{\mathrm{cm},b} + \tfrac{m_e}{M_{\mathrm{tot}}}\textstyle\sum_l r'_{l,b}\bigr\} \nonumber \\
&= \underbrace{\{p'_{i,a},R_{\mathrm{cm},b}\}}_{0} + \tfrac{m_e}{M_{\mathrm{tot}}}\!\sum_l\!\underbrace{\{p'_{i,a},r'_{l,b}\}}_{-\delta_{il}\delta_{ab}} \nonumber \\
&\qquad - \tfrac{m_e}{M_{\mathrm{tot}}}\!\underbrace{\{P_{\mathrm{cm},a},R_{\mathrm{cm},b}\}}_{-\delta_{ab}} - \tfrac{m_e^2}{M_{\mathrm{tot}}^2}\sum_l\!\underbrace{\{P_{\mathrm{cm},a},r'_{l,b}\}}_{0} \nonumber\\
&= -\tfrac{m_e}{M_{\mathrm{tot}}}\delta_{ab} + \tfrac{m_e}{M_{\mathrm{tot}}}\delta_{ab} = 0. \tag{2.16}
\end{align}
All other elementary brackets are preserved by inspection (each is either unchanged or trivially zero), so (2.15) is canonical.

\subsubsection{Transformed Hamiltonian}

Invert (2.15) for substitution into (2.12): $\mathbf{p}'_i = \tilde{\mathbf{p}}'_i + (m_e/M_{\mathrm{tot}})\mathbf{P}_{\mathrm{cm}}$. The potential $V$ depends only on $\mathbf{R}'_j$ and $\mathbf{r}'_i$, which are unchanged, so $V$ is form-invariant. Working term by term.

\noindent\emph{Cross-coupling (B):}
\begin{align}
-\frac{\mathbf{P}_{\mathrm{cm}}\!\cdot\!\sum_i\mathbf{p}'_i}{M_N}
&= -\frac{\mathbf{P}_{\mathrm{cm}}\!\cdot\!\bigl(\sum_i\tilde{\mathbf{p}}'_i + N_{\mathrm{elec}}\tfrac{m_e}{M_{\mathrm{tot}}}\mathbf{P}_{\mathrm{cm}}\bigr)}{M_N} \nonumber\\
&= -\frac{\mathbf{P}_{\mathrm{cm}}\!\cdot\!\sum_i\tilde{\mathbf{p}}'_i}{M_N} \;-\; \frac{N_{\mathrm{elec}} m_e\,\mathbf{P}_{\mathrm{cm}}^{\,2}}{M_N M_{\mathrm{tot}}}. \tag{2.17}
\end{align}

\noindent\emph{Mass-polarization quadratic:}
\begin{align}
\frac{(\sum_i\mathbf{p}'_i)^{2}}{2M_N}
&= \frac{\bigl(\sum_i\tilde{\mathbf{p}}'_i + N_{\mathrm{elec}}\tfrac{m_e}{M_{\mathrm{tot}}}\mathbf{P}_{\mathrm{cm}}\bigr)^{2}}{2M_N} \nonumber\\
&= \frac{(\sum_i\tilde{\mathbf{p}}'_i)^{2}}{2M_N} \;+\; \frac{N_{\mathrm{elec}} m_e\,\mathbf{P}_{\mathrm{cm}}\!\cdot\!\sum_i\tilde{\mathbf{p}}'_i}{M_N M_{\mathrm{tot}}} \;+\; \frac{N_{\mathrm{elec}}^2 m_e^2\,\mathbf{P}_{\mathrm{cm}}^{\,2}}{2 M_N M_{\mathrm{tot}}^2}. \tag{2.18}
\end{align}

\noindent\emph{Electron kinetic energy:}
\begin{align}
\sum_i\frac{\mathbf{p}_i^{'\,2}}{2m_e}
&= \sum_i\frac{(\tilde{\mathbf{p}}'_i + \tfrac{m_e}{M_{\mathrm{tot}}}\mathbf{P}_{\mathrm{cm}})^{2}}{2m_e} \nonumber\\
&= \sum_i\frac{\tilde{\mathbf{p}}_i^{'\,2}}{2m_e} \;+\; \frac{\mathbf{P}_{\mathrm{cm}}\!\cdot\!\sum_i\tilde{\mathbf{p}}'_i}{M_{\mathrm{tot}}} \;+\; \frac{N_{\mathrm{elec}} m_e\,\mathbf{P}_{\mathrm{cm}}^{\,2}}{2 M_{\mathrm{tot}}^2}. \tag{2.19}
\end{align}

\noindent\emph{Coefficient of $\mathbf{P}_{\mathrm{cm}}\!\cdot\!\sum_i\tilde{\mathbf{p}}'_i$:} Summing (2.17), (2.18), (2.19),
\begin{equation}
-\frac{1}{M_N} \;+\; \frac{N_{\mathrm{elec}} m_e}{M_N M_{\mathrm{tot}}} \;+\; \frac{1}{M_{\mathrm{tot}}} \;=\; \frac{-M_{\mathrm{tot}} + N_{\mathrm{elec}} m_e + M_N}{M_N M_{\mathrm{tot}}} \;=\; 0, \tag{2.20}
\end{equation}
using $M_{\mathrm{tot}} = M_N + N_{\mathrm{elec}} m_e$. \emph{The cross-coupling (B) is killed identically.}

\noindent\emph{Coefficient of $\mathbf{P}_{\mathrm{cm}}^{\,2}$:} Adding the $\mathbf{P}_{\mathrm{cm}}^{\,2}/(2M_N)$ from (A) plus the residual $\mathbf{P}_{\mathrm{cm}}^{\,2}$-pieces from (2.17)--(2.19),
\begin{align*}
\frac{1}{2M_N} - \frac{N_{\mathrm{elec}} m_e}{M_N M_{\mathrm{tot}}} + \frac{N_{\mathrm{elec}}^2 m_e^2}{2 M_N M_{\mathrm{tot}}^2} + \frac{N_{\mathrm{elec}} m_e}{2 M_{\mathrm{tot}}^2}
&= \frac{1}{2M_N} - \frac{N_{\mathrm{elec}} m_e}{M_N M_{\mathrm{tot}}} + \frac{N_{\mathrm{elec}} m_e}{2M_{\mathrm{tot}}^2}\!\left(\frac{N_{\mathrm{elec}} m_e}{M_N} + 1\right) \nonumber\\
&= \frac{1}{2M_N} - \frac{N_{\mathrm{elec}} m_e}{M_N M_{\mathrm{tot}}} + \frac{N_{\mathrm{elec}} m_e}{2M_{\mathrm{tot}}^2}\cdot\frac{M_{\mathrm{tot}}}{M_N} \nonumber\\
&= \frac{1}{2M_N} - \frac{N_{\mathrm{elec}} m_e}{M_N M_{\mathrm{tot}}} + \frac{N_{\mathrm{elec}} m_e}{2 M_N M_{\mathrm{tot}}} \nonumber\\
&= \frac{1}{2M_N} - \frac{N_{\mathrm{elec}} m_e}{2 M_N M_{\mathrm{tot}}} = \frac{M_{\mathrm{tot}}-N_{\mathrm{elec}} m_e}{2 M_N M_{\mathrm{tot}}} = \frac{M_N}{2 M_N M_{\mathrm{tot}}} = \frac{1}{2M_{\mathrm{tot}}}. \tag{2.21}
\end{align*}
\emph{The wrong-mass (A) is corrected to $M_{\mathrm{tot}}$.}

\noindent Assembling everything,
\begin{equation}
\boxed{\;\tilde H = \frac{\mathbf{P}_{\mathrm{cm}}^{\,2}}{2M_{\mathrm{tot}}} \;+\; \sum_j\frac{\mathbf{P}_j^{'\,2}}{2M_j} \;+\; \sum_i\frac{\tilde{\mathbf{p}}_i^{'\,2}}{2m_e} \;+\; \frac{1}{2M_N}\Bigl(\sum_i\tilde{\mathbf{p}}'_i\Bigr)^{\!2} \;+\; V(\mathbf{R}'_j,\mathbf{r}'_i).\;} \tag{2.22}
\end{equation}
The COM kinetic energy carries the \emph{correct} total mass $M_{\mathrm{tot}}$. The internal piece is structurally identical to (2.11): the same mass-polarization $(\sum_i\tilde{\mathbf{p}}'_i)^2/(2M_N)$, the same nuclear kinetic energy, the same potential. Since $\tilde H$ is independent of $\tilde{\mathbf{R}}_{\mathrm{cm}}$, $\tilde{\mathbf{P}}_{\mathrm{cm}}=\mathbf{P}_{\mathrm{cm}}$ is conserved, and one may proceed with the internal dynamics for \emph{any} value of $\mathbf{P}_{\mathrm{cm}}$ --- not just zero --- with no kinetic artifact.

\subsubsection{Geometric interpretation}

The new ``COM coordinate'' is
\begin{equation}
\tilde{\mathbf{R}}_{\mathrm{cm}} = \mathbf{R}_{\mathrm{cm}} + \frac{m_e}{M_{\mathrm{tot}}}\sum_i\mathbf{r}'_i = \frac{M_N\mathbf{R}_{\mathrm{cm}} + m_e\sum_i\mathbf{r}_i}{M_{\mathrm{tot}}} = \mathbf{R}_{\mathrm{tot}}^{\mathrm{CoM}}, \tag{2.23}
\end{equation}
using $\mathbf{r}_i = \mathbf{R}_{\mathrm{cm}} + \mathbf{r}'_i$. The canonical transformation has \emph{quietly relabeled} the COM coordinate from nuclear to total COM, while \emph{leaving the relative coordinates $\mathbf{R}'_j$ and $\mathbf{r}'_i$ untouched}. This is precisely what makes the framework still ``nuclear-COM-based'': the Born--Oppenheimer-friendly relative coordinates of Section 2.1 survive unchanged. The momentum shift $\tilde{\mathbf{p}}'_i = \mathbf{p}'_i - (m_e/M_{\mathrm{tot}})\mathbf{P}_{\mathrm{cm}}$ is the kinetic compensation required to keep the Poisson brackets canonical: in Galilean terms, each electron's lab-frame momentum is reduced by $m_e\mathbf{V}_{\mathrm{CoM}}$ with $\mathbf{V}_{\mathrm{CoM}} = \mathbf{P}_{\mathrm{cm}}/M_{\mathrm{tot}}$.

The canonical transformation has thus cleanly separated the total COM motion \emph{exactly and for any} $\mathbf{P}_{\mathrm{cm}}$: in (2.22) the COM appears only through the decoupled free term $\mathbf{P}_{\mathrm{cm}}^{\,2}/2M_{\mathrm{tot}}$ --- carrying the \emph{correct} total mass $M_{\mathrm{tot}}$ and free of the cross-coupling (B) --- while every other term depends only on the internal variables. It is this total Hamiltonian (2.22), exact for \emph{any} $\mathbf{P}_{\mathrm{cm}}$ and not merely the rest frame, that we carry into the body-fixed treatment of Section 3.

\section{BF frame for both nuclei and electrons}

\subsection{Setup}

With Eckart's frame~\cite{eckart1935} fixed by the \emph{nuclei},

\begin{equation}
\mathbf{R}_j'=\mathbf{S}(\phi,\theta,\chi)\,\bar{\mathbf{R}}_j(Q),\qquad \mathbf{r}_i'=\mathbf{S}(\phi,\theta,\chi)\,\bar{\mathbf{r}}_i. \tag{3.1}
\end{equation}
The Eckart conditions $\sum_j M_j\bar{\mathbf{R}}_j=0$ and $\sum_j M_j\bar{\mathbf{R}}_j^{\mathrm{eq}}\!\times\bar{\mathbf{R}}_j=0$ fix the orientation of $\mathbf{S}$ from the \emph{nuclear} positions; the same $\mathbf{S}$ is then used to define BF electron coordinates $\bar{\mathbf{r}}_i$. The nuclear shape is parametrized by rectilinear mass-weighted \textbf{normal coordinates} $Q_b$,
\begin{equation}
\bar{\mathbf{R}}_j(Q) = \bar{\mathbf{R}}_j^{\mathrm{eq}} + \frac{1}{\sqrt{M_j}}\sum_{b}\mathbf{l}_{jb}\,Q_b , \tag{3.1a}
\end{equation}
with mass-weighted normal-mode vectors $\mathbf{l}_{jb}$ obeying orthonormality and the Eckart (Sayvetz) conditions~\cite{eckart1935,sayvetz1939}
\begin{equation}
\sum_j \mathbf{l}_{ja}\!\cdot\!\mathbf{l}_{jb}=\delta_{ab},\qquad
\sum_j\sqrt{M_j}\,\mathbf{l}_{jb}=0,\qquad
\sum_j\sqrt{M_j}\,\bar{\mathbf{R}}_j^{\mathrm{eq}}\!\times\!\mathbf{l}_{jb}=0 , \tag{3.1b}
\end{equation}
so that $\partial\bar{\mathbf{R}}_j/\partial Q_b=\mathbf{l}_{jb}/\sqrt{M_j}$ is constant. \\
BF angular velocity:
\begin{equation}
(\mathbf{S}^T\dot{\mathbf{S}})_{\alpha\beta}=-\epsilon_{\alpha\beta\gamma}\,\omega_\gamma \Longrightarrow \mathbf{S}^T\dot{\mathbf{S}}\,\mathbf{v}=\boldsymbol{\omega}\times\mathbf{v}. \tag{3.2}
\end{equation}
Velocities (using $|\mathbf{S}\mathbf{v}|^2=|\mathbf{v}|^2$):
\begin{equation}
\dot{\mathbf{R}}_j'=\mathbf{S}\bigl(\boldsymbol{\omega}\!\times\!\bar{\mathbf{R}}_j+\dot{\bar{\mathbf{R}}}_j\bigr),\qquad \dot{\mathbf{r}}_i'=\mathbf{S}\bigl(\boldsymbol{\omega}\!\times\!\bar{\mathbf{r}}_i+\dot{\bar{\mathbf{r}}}_i\bigr). \tag{3.3}
\end{equation}

\subsection{From the internal Hamiltonian to a working Routhian}

I need an action principle for the body-fixed \emph{nuclear} coordinates. The electrons, however, are best left in the canonical (momentum) form they already carry in (2.22): their kinetic energy and mass polarization form a rotational scalar that passes through the body-fixed transformation unchanged. I therefore Legendre-transform \emph{only} the COM and nuclear momenta back to velocities --- keeping $\mathbf{P}_{\mathrm{cm}}$ arbitrary, never setting it to zero --- and keep the electron momenta $\tilde{\mathbf{p}}'_i$ as canonical variables. The relevant Hamilton equations from (2.22) are
\begin{equation}
\dot{\tilde{\mathbf{R}}}_{\mathrm{cm}}=\mathbf{P}_{\mathrm{cm}}/M_{\mathrm{tot}},\qquad \dot{\mathbf{R}}_j'=\mathbf{P}'_j/M_j,\qquad \dot{\mathbf{r}}_i'=\tilde{\mathbf{p}}_i'/m_e+(\textstyle\sum_l\tilde{\mathbf{p}}_l')/M_N, \tag{3.4}
\end{equation}
the last of which inverts to the electron momentum--velocity relation [recorded for the velocity-form cross-check below, around (3.18)]
\begin{equation}
\tilde{\mathbf{p}}_i'=m_e\dot{\mathbf{r}}_i'-\frac{m_e^2}{M_{\mathrm{tot}}}\sum_l\dot{\mathbf{r}}_l',\qquad \sum_l\tilde{\mathbf{p}}_l'=\frac{m_e M_N}{M_{\mathrm{tot}}}\,\sum_i\dot{\mathbf{r}}_i'. \tag{3.5}
\end{equation}
Legendre-transforming the COM and nuclear momenta only, $\mathcal{R}=\mathbf{P}_{\mathrm{cm}}\!\cdot\!\dot{\tilde{\mathbf{R}}}_{\mathrm{cm}}+\sum_j\mathbf{P}'_j\!\cdot\!\dot{\mathbf{R}}'_j-\tilde H$, yields the \textbf{Routhian}
\begin{equation}
\boxed{\;\mathcal{R}=\tfrac12 M_{\mathrm{tot}}|\dot{\tilde{\mathbf{R}}}_{\mathrm{cm}}|^2+\tfrac12\sum_j M_j|\dot{\mathbf{R}}_j'|^2-H_e(\tilde{\mathbf{p}}'_i)-V,\qquad H_e\equiv\sum_i\frac{\tilde{\mathbf{p}}_i'^{\,2}}{2m_e}+\frac{1}{2M_N}\Bigl(\sum_i\tilde{\mathbf{p}}'_i\Bigr)^{2},\;} \tag{3.6}
\end{equation}
which is a Lagrangian in $(\tilde{\mathbf{R}}_{\mathrm{cm}},\mathbf{R}'_j)$ and \emph{minus} a Hamiltonian in $(\mathbf{r}'_i,\tilde{\mathbf{p}}'_i)$: extremizing $\int\mathcal{R}\,dt$ gives Euler--Lagrange equations for the COM and nuclei and Hamilton's equations for the electrons. The electronic kinetic energy and mass polarization thus retain their (2.22) momentum form $H_e$ --- they are never converted to velocity form. (The overall minus sign on $H_e$ is the Routhian's partial-Legendre signature; the mass polarization keeps its physical positive coefficient $+1/2M_N$ inside $H_e$, and the whole $H_e$ re-enters the Hamiltonian with a plus sign when it is reassembled in (3.22).)

\subsection{Body-fixed Routhian}

Substitute the body-fixed velocities (3.3) into the nuclear term of (3.6); the COM term is untouched by the rotation. The electron Hamiltonian $H_e$ is a \emph{rotational scalar} --- it is built from the dot products $\tilde{\mathbf{p}}'_i\!\cdot\!\tilde{\mathbf{p}}'_l$ --- so under the passive rotation to the body frame it is unchanged in form. Defining the body-fixed electron momenta $\bar{\mathbf{p}}_i\equiv\mathbf{S}^{T}\tilde{\mathbf{p}}'_i$ (a proper rotation, so $|\bar{\mathbf{p}}_i|^2=|\tilde{\mathbf{p}}'_i|^2$ and $\sum_i\bar{\mathbf{p}}_i=\mathbf{S}^{T}\sum_i\tilde{\mathbf{p}}'_i$, whence $(\sum_i\bar{\mathbf{p}}_i)^2=(\sum_i\tilde{\mathbf{p}}'_i)^2$),
\begin{equation}
H_e=\sum_i\frac{\tilde{\mathbf{p}}_i'^{\,2}}{2m_e}+\frac{1}{2M_N}\Bigl(\sum_i\tilde{\mathbf{p}}'_i\Bigr)^{2}=\sum_i\frac{\bar{\mathbf{p}}_i^{\,2}}{2m_e}+\frac{1}{2M_N}\Bigl(\sum_i\bar{\mathbf{p}}_i\Bigr)^{2}. \tag{3.6a}
\end{equation}
Defining $\bar{\mathbf{r}}\equiv\sum_i\bar{\mathbf{r}}_i$, the body-fixed Routhian is therefore
\begin{equation}
\mathcal{R}=\tfrac12 M_{\mathrm{tot}}|\dot{\tilde{\mathbf{R}}}_{\mathrm{cm}}|^2+\tfrac12\!\sum_j M_j|\boldsymbol{\omega}\!\times\!\bar{\mathbf{R}}_j+\dot{\bar{\mathbf{R}}}_j|^2-\sum_i\frac{\bar{\mathbf{p}}_i^{\,2}}{2m_e}-\frac{1}{2M_N}\Bigl(\sum_i\bar{\mathbf{p}}_i\Bigr)^{2}-V, \tag{3.7}
\end{equation}
in which only the nuclear kinetic term carries the body-fixed angular velocity $\boldsymbol{\omega}$. The potential is now manifestly $\mathbf{S}$-independent:
\begin{equation}
V=\frac{1}{4\pi\epsilon_0}\left(\frac{1}{2}\!\sum_{j\ne k}\!\frac{Z_jZ_ke^2}{\bar R_{jk}}-\!\sum_{j,i}\!\frac{Z_je^2}{|\bar{\mathbf{R}}_j-\bar{\mathbf{r}}_i|}+\frac{1}{2}\!\sum_{i\ne l}\!\frac{e^2}{\bar{r}_{il}}\right). \tag{3.8}
\end{equation}
Rotational symmetry is rendered explicit; the BF electronic problem at fixed $Q$ is self-contained.

\subsection{Expanding (3.7)}

\subsubsection{Generic squared-term identity}

For any vectors $\mathbf{a},\dot{\mathbf{a}},\boldsymbol{\omega}\in\mathbb{R}^3$, the BCH-style identity $\mathbf{a}\!\times\!(\boldsymbol{\omega}\!\times\!\mathbf{a})=\boldsymbol{\omega}\,a^2-\mathbf{a}(\mathbf{a}\!\cdot\!\boldsymbol{\omega})$ gives
\begin{align}
|\boldsymbol{\omega}\!\times\!\mathbf{a}|^{2} 
&= (\boldsymbol{\omega}\!\times\!\mathbf{a})\!\cdot\!(\boldsymbol{\omega}\!\times\!\mathbf{a})
\;=\; \boldsymbol{\omega}\!\cdot\!\bigl[\mathbf{a}\!\times\!(\boldsymbol{\omega}\!\times\!\mathbf{a})\bigr] \nonumber\\
&= \boldsymbol{\omega}\!\cdot\!\bigl[\boldsymbol{\omega}\,a^{2}-\mathbf{a}(\mathbf{a}\!\cdot\!\boldsymbol{\omega})\bigr]
\;=\; \boldsymbol{\omega}^{T}(a^{2}\mathbb{1}-\mathbf{a}\mathbf{a}^{T})\,\boldsymbol{\omega}, \tag{3.9a}
\end{align}
while the scalar triple product gives
\begin{equation}
(\boldsymbol{\omega}\!\times\!\mathbf{a})\!\cdot\!\dot{\mathbf{a}}
\;=\; \boldsymbol{\omega}\!\cdot\!(\mathbf{a}\!\times\!\dot{\mathbf{a}}). \tag{3.9b}
\end{equation}
Hence, for any constant mass-like coefficient $m$,
\begin{align}
\tfrac{m}{2}\bigl|\boldsymbol{\omega}\!\times\!\mathbf{a} + \dot{\mathbf{a}}\bigr|^{2}
&= \tfrac{m}{2}\bigl[|\boldsymbol{\omega}\!\times\!\mathbf{a}|^{2} + 2(\boldsymbol{\omega}\!\times\!\mathbf{a})\!\cdot\!\dot{\mathbf{a}} + |\dot{\mathbf{a}}|^{2}\bigr] \nonumber\\
&= \underbrace{\tfrac{m}{2}\,\boldsymbol{\omega}^{T}(a^{2}\mathbb{1}-\mathbf{a}\mathbf{a}^{T})\boldsymbol{\omega}}_{\text{rotational}} \;+\; \underbrace{m\,\boldsymbol{\omega}\!\cdot\!(\mathbf{a}\!\times\!\dot{\mathbf{a}})}_{\text{Coriolis-type}} \;+\; \underbrace{\tfrac{m}{2}|\dot{\mathbf{a}}|^{2}}_{\text{vibrational}}. \tag{3.9c}
\end{align}

\subsubsection{Application to the nuclear term}

Only the nuclear kinetic term of the Routhian (3.7) carries $\boldsymbol{\omega}$. Apply (3.9c) with $m\!\to\! M_j$, $\mathbf{a}\!\to\!\bar{\mathbf{R}}_j$ and sum over $j$:
\begin{equation}
\tfrac12\!\sum_j M_j\bigl|\boldsymbol{\omega}\!\times\!\bar{\mathbf{R}}_j+\dot{\bar{\mathbf{R}}}_j\bigr|^{2}
\;=\; \tfrac12\boldsymbol{\omega}^{T}\!\mathbf{I}_N\boldsymbol{\omega} \;+\; \boldsymbol{\omega}\!\cdot\!\mathbf{L}_N^{\mathrm{vib}} \;+\; T_{\mathrm{vib}}^{N}. \tag{3.9d}
\end{equation}
The electronic kinetic energy is \emph{not} expanded this way --- it is carried as the rotational scalar $H_e$ of (3.6a). Only for the velocity-form cross-check of \S\,3.6 do we record the decomposition that the electronic and mass-polarization terms \emph{would} produce if the electrons were instead kept in velocity form,
\begin{equation}
\tfrac{m_e}{2}\!\sum_i\bigl|\boldsymbol{\omega}\!\times\!\bar{\mathbf{r}}_i+\dot{\bar{\mathbf{r}}}_i\bigr|^{2}
= \tfrac12\boldsymbol{\omega}^{T}\!\mathbf{I}_e\boldsymbol{\omega} + \boldsymbol{\omega}\!\cdot\!\mathbf{L}_e + T_{e}^{\mathrm{BF}},\qquad
-\tfrac{m_e^{2}}{2M_{\mathrm{tot}}}\bigl|\boldsymbol{\omega}\!\times\!\bar{\mathbf{r}}+\dot{\bar{\mathbf{r}}}\bigr|^{2}
= -\tfrac12\boldsymbol{\omega}^{T}\!\mathbf{I}_X\boldsymbol{\omega} - \boldsymbol{\omega}\!\cdot\!\mathbf{L}_X - T_{X}^{\mathrm{mp}}. \tag{3.9e}
\end{equation}

The constituent inertia tensors are
\begin{equation}
\mathbf{I}_N=\sum_j M_j(\bar R_j^2\mathbb{1}-\bar{\mathbf{R}}_j\bar{\mathbf{R}}_j^T),\qquad \mathbf{I}_e=m_e\!\sum_i(\bar r_i^2\mathbb{1}-\bar{\mathbf{r}}_i\bar{\mathbf{r}}_i^T), \tag{3.9}
\end{equation}
\begin{equation}
\mathbf{I}_X=\tfrac{m_e^2}{M_{\mathrm{tot}}}(\bar r^2\mathbb{1}-\bar{\mathbf{r}}\bar{\mathbf{r}}^T), \tag{3.10}
\end{equation}
the angular-momentum-like (velocity-form) quantities are
\begin{equation}
\mathbf{L}_N^{\mathrm{vib}}=\sum_jM_j\bar{\mathbf{R}}_j\!\times\!\dot{\bar{\mathbf{R}}}_j,\quad \mathbf{L}_e=m_e\!\sum_i\bar{\mathbf{r}}_i\!\times\!\dot{\bar{\mathbf{r}}}_i,\quad \mathbf{L}_X=\tfrac{m_e^2}{M_{\mathrm{tot}}}\bar{\mathbf{r}}\!\times\!\dot{\bar{\mathbf{r}}}, \tag{3.11}
\end{equation}
and the ``vibrational'' kinetic pieces are
\begin{equation}
T_{\mathrm{vib}}^{N}=\tfrac12\!\sum_jM_j|\dot{\bar{\mathbf{R}}}_j|^2=\tfrac12\!\sum_b\dot Q_b^2,\quad T_{e}^{\mathrm{BF}}=\tfrac{m_e}{2}\!\sum_i|\dot{\bar{\mathbf{r}}}_i|^2,\quad T_X^{\mathrm{mp}}=\tfrac{m_e^2}{2M_{\mathrm{tot}}}|\dot{\bar{\mathbf{r}}}|^2. \tag{3.12}
\end{equation}

\subsubsection{Assembled Routhian}

Inserting the nuclear expansion (3.9d) into (3.7), together with the decoupled COM term, the electron scalar $H_e$, and the potential:
\begin{equation}
\boxed{\;\mathcal{R}=\tfrac12 M_{\mathrm{tot}}|\dot{\tilde{\mathbf{R}}}_{\mathrm{cm}}|^2+\underbrace{\tfrac12\boldsymbol{\omega}^{T}\mathbf{I}_N\,\boldsymbol{\omega}+\boldsymbol{\omega}\!\cdot\!\mathbf{L}_N^{\mathrm{vib}}+T_{\mathrm{vib}}^{N}}_{\text{nuclear, velocity form}}\;-\;H_e\;-\;V.\;} \tag{3.13}
\end{equation}
The rotational kernel is the \textbf{nuclear} inertia $\mathbf{I}_N$ \emph{alone}: because the electrons are carried as the scalar $H_e$ rather than expanded in $\boldsymbol{\omega}$, no electronic or mass-polarization inertia ($\mathbf{I}_e$, $\mathbf{I}_X$) enters the kernel. The electrons couple to rotation only \emph{kinematically}, through their angular momentum, which we identify next. The body-fixed electron momenta and the \textbf{electronic angular momentum} are
\begin{equation}
\bar{\mathbf{p}}_i=\mathbf{S}^{T}\tilde{\mathbf{p}}'_i,\qquad \mathbf{L}_{\mathrm{elec}}\equiv\sum_i\bar{\mathbf{r}}_i\!\times\!\bar{\mathbf{p}}_i. \tag{3.14}
\end{equation}
\textbf{Remark.} The electronic and mass-polarization inertia/angular-momentum quantities $\mathbf{I}_e,\mathbf{I}_X,\mathbf{L}_e,\mathbf{L}_X$ of (3.9)--(3.11), and the velocity-form decomposition (3.9e), are needed \emph{only} in the cross-check of \S\,3.6; they play no role in the Routhian (3.13) itself.

\subsection{Canonical momenta and the total angular momentum}

Compute the conjugate momenta from the Routhian $\mathcal{R}$. The cyclic COM coordinate gives the conserved total momentum $\mathbf{P}_{\mathrm{cm}}=\partial \mathcal{R}/\partial\dot{\tilde{\mathbf{R}}}_{\mathrm{cm}}=M_{\mathrm{tot}}\dot{\tilde{\mathbf{R}}}_{\mathrm{cm}}$; being decoupled, it contributes nothing to the internal momenta below. The electron momenta $\bar{\mathbf{p}}_i$ are already canonical [eq.~(3.14)]; it remains to find the rotational and vibrational momenta, most transparently by differentiating the unexpanded nuclear term of (3.7).

\subsubsection{Nuclear angular momentum $\mathbf{J}_N=\partial \mathcal{R}/\partial\boldsymbol{\omega}$}

In the Routhian (3.13) only the nuclear rotational kernel and the nuclear Coriolis piece carry $\boldsymbol{\omega}$ --- the electron scalar $H_e$ and the COM term do not. The momentum conjugate to the body-fixed angular velocity is therefore the \emph{nuclear} angular momentum
\begin{equation}
\mathbf{J}_N\;\equiv\;\frac{\partial \mathcal{R}}{\partial\boldsymbol{\omega}}=\mathbf{I}_N\,\boldsymbol{\omega}+\mathbf{L}_N^{\mathrm{vib}}. \tag{3.15}
\end{equation}

\subsubsection{Total angular momentum $\mathbf{J}=\mathbf{J}_N+\mathbf{L}_{\mathrm{elec}}$}

The body-fixed frame $\mathbf{S}(\Omega)$ rotates the nuclei \emph{and} the electrons together [eq.~(3.1)], so the total angular momentum conjugate to the Euler angles --- the physical, conserved rovibronic angular momentum --- is the sum of the nuclear and electronic contributions. The electronic part is the canonical electron angular momentum $\mathbf{L}_{\mathrm{elec}}=\sum_i\bar{\mathbf{r}}_i\!\times\!\bar{\mathbf{p}}_i$ of (3.14), the body-fixed image of $\sum_i\mathbf{r}'_i\!\times\!\tilde{\mathbf{p}}'_i$ (the cross product co-rotates with $\mathbf{S}$). Hence
\begin{equation}
\mathbf{J}=\mathbf{J}_N+\mathbf{L}_{\mathrm{elec}}=\mathbf{I}_N\boldsymbol{\omega}+\mathbf{L}_N^{\mathrm{vib}}+\mathbf{L}_{\mathrm{elec}}. \tag{3.16}
\end{equation}

\subsubsection{Vibrational momentum $P_b=\partial \mathcal{R}/\partial\dot Q_b$}

Only $\dot{\bar{\mathbf{R}}}_j$ carries $\dot Q_b$ dependence; with $\partial\bar{\mathbf{R}}_j/\partial Q_b=\mathbf{l}_{jb}/\sqrt{M_j}$ from (3.1a),
\begin{equation}
P_b\;\equiv\;\frac{\partial \mathcal{R}}{\partial\dot Q_b}=\sum_jM_j\bigl(\boldsymbol{\omega}\!\times\!\bar{\mathbf{R}}_j+\dot{\bar{\mathbf{R}}}_j\bigr)\!\cdot\!\frac{\partial\bar{\mathbf{R}}_j}{\partial Q_b}=\sum_j\sqrt{M_j}\,\bigl(\boldsymbol{\omega}\!\times\!\bar{\mathbf{R}}_j+\dot{\bar{\mathbf{R}}}_j\bigr)\!\cdot\!\mathbf{l}_{jb}. \tag{3.17}
\end{equation}
Separating the rotational and vibrational pieces defines the \textbf{Coriolis vectors} $\boldsymbol{\sigma}_b$,
\begin{equation}
\boldsymbol{\sigma}_b \equiv \sum_jM_j\,\bar{\mathbf{R}}_j\!\times\!\frac{\partial\bar{\mathbf{R}}_j}{\partial Q_b}=\sum_j\sqrt{M_j}\,\bar{\mathbf{R}}_j\!\times\!\mathbf{l}_{jb}=\sum_a\boldsymbol{\zeta}_{ab}\,Q_a,\qquad \boldsymbol{\zeta}_{ab}\equiv\sum_j\mathbf{l}_{ja}\!\times\!\mathbf{l}_{jb}, \tag{3.17a}
\end{equation}
the second form following from (3.1a) with the Eckart condition $\sum_j\sqrt{M_j}\bar{\mathbf{R}}_j^{\mathrm{eq}}\!\times\!\mathbf{l}_{jb}=0$. The vibrational metric is the identity, $\sum_jM_j(\partial\bar{\mathbf{R}}_j/\partial Q_a)\!\cdot\!(\partial\bar{\mathbf{R}}_j/\partial Q_b)=\sum_j\mathbf{l}_{ja}\!\cdot\!\mathbf{l}_{jb}=\delta_{ab}$, so the canonical vibrational momentum carries \emph{no} Wilson $G$-matrix,
\begin{equation}
P_b \;=\; \boldsymbol{\omega}\!\cdot\!\boldsymbol{\sigma}_b(Q) \;+\; \dot Q_b. \tag{3.17b}
\end{equation}

\subsection{Key identity}

\subsubsection{The clean split}

The total-angular-momentum decomposition (3.16) \emph{is} the key identity: subtracting the electronic part from both sides,
\begin{equation}
\boxed{\;\mathbf{J}-\mathbf{L}_{\mathrm{elec}}=\mathbf{I}_N(Q)\,\boldsymbol{\omega}+\mathbf{L}_N^{\mathrm{vib}}.\;} \tag{3.19}
\end{equation}
The right side is a purely nuclear-velocity quantity, with the \emph{nuclear} inertia $\mathbf{I}_N$ as rotational kernel. The entire electronic and mass-polarization content sits on the left, in the kinematic shift $\mathbf{J}\to\mathbf{J}-\mathbf{L}_{\mathrm{elec}}$; it never enters the kernel, because in the Routhian (3.13) the electrons were carried as the scalar $H_e$ and contributed no inertia at all.

\subsubsection{Cross-check: the $\mathbf{I}_e/\mathbf{I}_X$ cancellation of the velocity-form route}

Had we instead converted the electrons to velocity form --- expanding their kinetic energy via (3.9e) --- the electronic inertia and Coriolis terms would have entered \emph{both} $\mathbf{J}$ and $\mathbf{L}_{\mathrm{elec}}$, and (3.19) would emerge only after an exact cancellation. The velocity-form electron momentum is $\bar{\mathbf{p}}_i=m_e(\boldsymbol{\omega}\!\times\!\bar{\mathbf{r}}_i+\dot{\bar{\mathbf{r}}}_i)-\tfrac{m_e^2}{M_{\mathrm{tot}}}(\boldsymbol{\omega}\!\times\!\bar{\mathbf{r}}+\dot{\bar{\mathbf{r}}})$ [the body-fixed image of (3.5)]; substituting into $\mathbf{L}_{\mathrm{elec}}=\sum_i\bar{\mathbf{r}}_i\!\times\!\bar{\mathbf{p}}_i$ and using the BAC--CAB identity $\mathbf{a}\!\times\!(\boldsymbol{\omega}\!\times\!\mathbf{a})=(a^2\mathbb{1}-\mathbf{a}\mathbf{a}^{T})\boldsymbol{\omega}$ on the two $\boldsymbol{\omega}$-terms collapses the four pieces to
\begin{equation}
\mathbf{L}_{\mathrm{elec}}=(\mathbf{I}_e-\mathbf{I}_X)\boldsymbol{\omega}+(\mathbf{L}_e-\mathbf{L}_X), \tag{3.18}
\end{equation}
with $\mathbf{I}_e,\mathbf{I}_X,\mathbf{L}_e,\mathbf{L}_X$ from (3.9)--(3.11). The velocity-form total angular momentum would read $\mathbf{J}=(\mathbf{I}_N+\mathbf{I}_e-\mathbf{I}_X)\boldsymbol{\omega}+(\mathbf{L}_N^{\mathrm{vib}}+\mathbf{L}_e-\mathbf{L}_X)$ [from (3.9d)+(3.9e)], so
\begin{equation*}
\mathbf{J}-\mathbf{L}_{\mathrm{elec}}=\mathbf{I}_N\boldsymbol{\omega}+\mathbf{L}_N^{\mathrm{vib}}+\underbrace{\bigl[(\mathbf{I}_e-\mathbf{I}_X)\boldsymbol{\omega}+(\mathbf{L}_e-\mathbf{L}_X)\bigr]-\mathbf{L}_{\mathrm{elec}}}_{=\,0\ \text{by (3.18)}}=\mathbf{I}_N\boldsymbol{\omega}+\mathbf{L}_N^{\mathrm{vib}},
\end{equation*}
reproducing (3.19): the electronic inertia $\mathbf{I}_e$ and the mass-polarization inertia $\mathbf{I}_X$ cancel \emph{exactly}. This ``algebraic miracle'' of the Eckart construction is automatic in the mixed treatment, where $\mathbf{I}_e$ and $\mathbf{I}_X$ never arise.

\subsubsection{Inversion: the Watson reciprocal inertia}

To invert (3.19) for $\boldsymbol{\omega}$ we must re-express the velocity-form $\mathbf{L}_N^{\mathrm{vib}}$ in canonical variables. In normal coordinates $\mathbf{L}_N^{\mathrm{vib}}=\sum_jM_j\bar{\mathbf{R}}_j\!\times\!\dot{\bar{\mathbf{R}}}_j=\sum_b\boldsymbol{\sigma}_b\,\dot Q_b$; eliminating $\dot Q_b=P_b-\boldsymbol{\omega}\!\cdot\!\boldsymbol{\sigma}_b$ from (3.17b),
\begin{equation}
\mathbf{L}_N^{\mathrm{vib}}=\sum_b\boldsymbol{\sigma}_b\bigl(P_b-\boldsymbol{\omega}\!\cdot\!\boldsymbol{\sigma}_b\bigr)=\boldsymbol{\pi}-\boldsymbol{\sigma}\boldsymbol{\sigma}^{\mathsf T}\boldsymbol{\omega},\qquad
\boldsymbol{\pi}\equiv\sum_b\boldsymbol{\sigma}_bP_b,\quad \boldsymbol{\sigma}\boldsymbol{\sigma}^{\mathsf T}\equiv\sum_b\boldsymbol{\sigma}_b\boldsymbol{\sigma}_b^{\mathsf T}, \tag{3.20a}
\end{equation}
with $\boldsymbol{\pi}$ the \textbf{field-free vibrational angular momentum}. Substituting into (3.19), $\mathbf{J}-\mathbf{L}_{\mathrm{elec}}=(\mathbf{I}_N-\boldsymbol{\sigma}\boldsymbol{\sigma}^{\mathsf T})\boldsymbol{\omega}+\boldsymbol{\pi}\equiv\mathbf{I}'\boldsymbol{\omega}+\boldsymbol{\pi}$, so
\begin{equation}
\boldsymbol{\omega}=\boldsymbol{\mu}\bigl(\mathbf{J}-\mathbf{L}_{\mathrm{elec}}-\boldsymbol{\pi}\bigr),\qquad
\boxed{\;\boldsymbol{\mu}\equiv(\mathbf{I}')^{-1}=\bigl(\mathbf{I}_N-\boldsymbol{\sigma}\boldsymbol{\sigma}^{\mathsf T}\bigr)^{-1},\;} \tag{3.20}
\end{equation}
the \textbf{Watson reciprocal-inertia tensor}~\cite{watson1968}. The rotation--vibration coupling is now carried by the modified inertia $\mathbf{I}'$ and the uniform shift $\mathbf{J}\to\mathbf{J}-\mathbf{L}_{\mathrm{elec}}-\boldsymbol{\pi}$ --- with no Wilson $G$-matrix.

\subsection{Hamiltonian in canonical variables}

We now assemble the body-fixed Hamiltonian in the canonical variables $(\Omega,\mathbf{J};\,Q,P_b;\,\bar{\mathbf{r}}_i,\bar{\mathbf{p}}_i)$. The total energy splits as $\tilde H=\mathbf{P}_{\mathrm{cm}}^2/2M_{\mathrm{tot}}+H_{\mathrm{int}}$, the COM piece being a decoupled spectator. The internal Hamiltonian is the sum of the nuclear rotation--vibration kinetic energy, the electron scalar, and the potential,
\begin{equation*}
H_{\mathrm{int}}=T_N^{\mathrm{rot+vib}}+H_e+V,\qquad T_N^{\mathrm{rot+vib}}\equiv\tfrac12\boldsymbol{\omega}^{T}\mathbf{I}_N\boldsymbol{\omega}+\boldsymbol{\omega}\!\cdot\!\mathbf{L}_N^{\mathrm{vib}}+T_{\mathrm{vib}}^{N},
\end{equation*}
read off the Routhian (3.13). The electron scalar $H_e$ [eq.~(3.6a)] and the potential $V$ are already canonical; only $T_N^{\mathrm{rot+vib}}$ must be re-expressed through the momenta $(\mathbf{J},P_b)$.

\subsubsection{Nuclear rotation--vibration block}

$T_N^{\mathrm{rot+vib}}$ is homogeneous-quadratic in the nuclear velocities $(\boldsymbol{\omega},\dot Q_b)$, so by Euler's theorem $2T_N^{\mathrm{rot+vib}}=\mathbf{J}_N\!\cdot\!\boldsymbol{\omega}+\sum_b P_b\dot Q_b$, with conjugate momenta $\mathbf{J}_N=\mathbf{J}-\mathbf{L}_{\mathrm{elec}}$ [eqs.~(3.15),(3.19)] and $P_b$ [eq.~(3.17b)]. Substituting $\dot Q_b=P_b-\boldsymbol{\omega}\!\cdot\!\boldsymbol{\sigma}_b$ and $\sum_b P_b\boldsymbol{\sigma}_b=\boldsymbol{\pi}$,
\begin{equation*}
2T_N^{\mathrm{rot+vib}}=(\mathbf{J}-\mathbf{L}_{\mathrm{elec}})\!\cdot\!\boldsymbol{\omega}+\sum_bP_b\bigl(P_b-\boldsymbol{\omega}\!\cdot\!\boldsymbol{\sigma}_b\bigr)=(\mathbf{J}-\mathbf{L}_{\mathrm{elec}}-\boldsymbol{\pi})\!\cdot\!\boldsymbol{\omega}+\sum_bP_b^2 .
\end{equation*}
Inserting $\boldsymbol{\omega}=\boldsymbol{\mu}(\mathbf{J}-\mathbf{L}_{\mathrm{elec}}-\boldsymbol{\pi})$ from (3.20),
\begin{equation}
\boxed{\;T_N^{\mathrm{rot+vib}}\Big|_{\mathrm{canonical}} \;=\; \tfrac12(\mathbf{J}-\mathbf{L}_{\mathrm{elec}}-\boldsymbol{\pi})^{T}\boldsymbol{\mu}\,(\mathbf{J}-\mathbf{L}_{\mathrm{elec}}-\boldsymbol{\pi}) + \tfrac12\sum_{b}P_b^2.\;} \tag{3.20d}
\end{equation}
The vibrational kinetic energy is the \emph{bare} $\tfrac12\sum_bP_b^2$; the entire rotation--vibration coupling is carried by $\boldsymbol{\pi}$ and by the Watson inertia $\boldsymbol{\mu}=(\mathbf{I}_N-\boldsymbol{\sigma}\boldsymbol{\sigma}^{\mathsf T})^{-1}$ --- no Wilson $G$-matrix survives.

\subsubsection{Electronic block}

The electronic kinetic energy and mass polarization require no further work: they are the rotational scalar $H_e$ of (3.6a), carried unchanged from the Step-2 Hamiltonian (2.22),
\begin{equation}
\boxed{\;H_e=\sum_i\frac{\bar{\mathbf{p}}_i^{\,2}}{2m_e}+\frac{1}{2M_N}\Bigl(\sum_i\bar{\mathbf{p}}_i\Bigr)^{2}.\;} \tag{3.21}
\end{equation}
Because $H_e$ is built from the rotation-invariant dot products $\bar{\mathbf{p}}_i\!\cdot\!\bar{\mathbf{p}}_l$ and $\bar{\mathbf{p}}_i=\mathbf{S}^{T}\tilde{\mathbf{p}}'_i$ is a passive rotation, its body-fixed form is \emph{identical} to the Step-2 form (2.22) under $\tilde{\mathbf{p}}'_i\to\bar{\mathbf{p}}_i$ --- no Legendre transform or angular-momentum expansion is needed. The mass polarization survives the body-fixed transformation intact in functional form.

\emph{Reduced-mass check} ($N_{\mathrm{elec}}=1$): $H_e=\bar{\mathbf{p}}_1^{\,2}/(2m_e)+\bar{\mathbf{p}}_1^{\,2}/(2M_N)=\bar{\mathbf{p}}_1^{\,2}/(2\mu_e)$ with $\mu_e=m_eM_N/(m_e+M_N)=m_eM_N/M_{\mathrm{tot}}$. $\checkmark$

\subsubsection{Final assembly: (3.22)}

Add (3.20d) and (3.21), include $V = V_{NN}(Q)+V_{Ne}(Q,\{\bar{\mathbf{r}}_i\})+V_{ee}(\{\bar{\mathbf{r}}_i\})$ from (3.8), and use $H_{\mathrm{int}} = T_N^{\mathrm{rot+vib}}+H_e+V$:
\begin{equation}
\boxed{
\begin{aligned}
H_{\mathrm{int}}\;=\;& \tfrac12\bigl(\mathbf{J}-\mathbf{L}_{\mathrm{elec}}-\boldsymbol{\pi}\bigr)^T\boldsymbol{\mu}(Q)\bigl(\mathbf{J}-\mathbf{L}_{\mathrm{elec}}-\boldsymbol{\pi}\bigr) \\
& +\tfrac12\!\sum_{b}P_b^2 \\
& +\sum_i\frac{\bar{\mathbf{p}}_i^{\,2}}{2m_e}+\frac{1}{2M_N}\!\Bigl(\sum_i\bar{\mathbf{p}}_i\Bigr)^{\!2} \\
& +V_{NN}(Q)+V_{Ne}(Q,\{\bar{\mathbf{r}}_i\})+V_{ee}(\{\bar{\mathbf{r}}_i\}).
\end{aligned}}\tag{3.22}
\end{equation}

Restoring the decoupled COM piece, the complete canonical Hamiltonian is $\tilde H = \mathbf{P}_{\mathrm{cm}}^2/2M_{\mathrm{tot}} + H_{\mathrm{int}}$ with $H_{\mathrm{int}}$ given by (3.22) --- exact for \emph{any} $\mathbf{P}_{\mathrm{cm}}$. The COM term is a conserved free-particle spectator that decouples from every body-fixed variable, so the quantization of Section 4 proceeds with $H_{\mathrm{int}}$.

Here $\boldsymbol{\pi}=\sum_b\boldsymbol{\sigma}_bP_b$ is the field-free vibrational angular momentum, with the Coriolis vectors $\boldsymbol{\sigma}_b=\sum_a\boldsymbol{\zeta}_{ab}Q_a$ of (3.17a) and constants $\boldsymbol{\zeta}_{ab}=\sum_j\mathbf{l}_{ja}\!\times\!\mathbf{l}_{jb}$; the vibrational metric being the identity, no Wilson $G$-matrix enters.

$\mathbf{L}_{\mathrm{elec}}=\sum_i\bar{\mathbf{r}}_i\!\times\!\bar{\mathbf{p}}_i$ is the BF canonical electronic angular momentum, and crucially:
\begin{itemize}
    \item $\boldsymbol{\mu}=(\mathbf{I}_N-\boldsymbol{\sigma}\boldsymbol{\sigma}^{\mathsf T})^{-1}$ is the \textbf{Watson reciprocal-inertia tensor}, built from nuclear quantities alone: \emph{no} electron or mass-polarization corrections enter the rotational kernel --- they were absorbed into $\mathbf{L}_{\mathrm{elec}}$ via (3.19).
    \item The electron contribution to rotation appears entirely as the Coriolis-like coupling $-\mathbf{J}^T\boldsymbol{\mu}\mathbf{L}_{\mathrm{elec}}$.
    \item The mass polarization is preserved in its $\Bigl(\sum_i\bar{\mathbf{p}}_i\Bigr)^{\!2}\!\!/(2M_N)$ form.
\end{itemize}

\subsection{Complete coupling enumeration}

Expanding the rotational kernel:
$$
\tfrac12\bigl(\mathbf{J}-\mathbf{L}_{\mathrm{elec}}-\boldsymbol{\pi}\bigr)^T\boldsymbol{\mu}\bigl(\mathbf{J}-\mathbf{L}_{\mathrm{elec}}-\boldsymbol{\pi}\bigr)=
$$
$$
\underbrace{\tfrac12\mathbf{J}^T\boldsymbol{\mu}\mathbf{J}}_{\text{rotation}}
\underbrace{-\mathbf{J}^T\boldsymbol{\mu}\boldsymbol{\pi}}_{\text{rot–vib Coriolis}}
\underbrace{-\mathbf{J}^T\boldsymbol{\mu}\mathbf{L}_{\mathrm{elec}}}_{\text{rot–elec Coriolis}}
\underbrace{+\boldsymbol{\pi}^{T}\boldsymbol{\mu}\mathbf{L}_{\mathrm{elec}}}_{\text{vib–elec angular}}
+\tfrac12\boldsymbol{\pi}^{T}\boldsymbol{\mu}\boldsymbol{\pi}+\tfrac12\mathbf{L}_{\mathrm{elec}}^{T}\boldsymbol{\mu}\mathbf{L}_{\mathrm{elec}}.
$$

All couplings in (3.22), referred to the standard analyses~\cite{watson1968,bunkerjensen}:
\begin{enumerate}
    \item \textbf{Pure rotation, anisotropic top:} $\tfrac12\mathbf{J}^T\boldsymbol{\mu}(Q)\mathbf{J}$ --- $Q$-dependent moments of inertia (rot–vib coupling through inertia).
    \item \textbf{Rot–vib Coriolis}~\cite{wilson1955}: $-\mathbf{J}^T\boldsymbol{\mu}\boldsymbol{\pi}$.
    \item \textbf{Rot–electron Coriolis:} $-\mathbf{J}^T\boldsymbol{\mu}\mathbf{L}_{\mathrm{elec}}$ --- the kinetic source of $\Lambda$-doubling~\cite{mullikenchristy1931,vanvleck1951}, Renner–Teller~\cite{renner1934}, and spin–rotation~\cite{curl1965} analogues.
    \item \textbf{Vib–electron angular coupling:} $\boldsymbol{\pi}^{T}\boldsymbol{\mu}\mathbf{L}_{\mathrm{elec}}$ --- exact non-adiabatic angular coupling intrinsic to BF embedding~\cite{bornoppenheimer1927,bornhuang1954,littlejohn1997}.
    \item \textbf{Quadratic vibrational angular momentum:} $\tfrac12\boldsymbol{\pi}^{T}\boldsymbol{\mu}\boldsymbol{\pi}$.
    \item \textbf{Quadratic electronic angular momentum:} $\tfrac12\mathbf{L}_{\mathrm{elec}}^{T}\boldsymbol{\mu}\mathbf{L}_{\mathrm{elec}}$.
    \item \textbf{Pure vibrational kinetic:} $\tfrac12\sum_{b}P_b^2$.
    \item \textbf{Electronic kinetic:} $\sum_i\bar{\mathbf{p}}_i^{\,2}/(2m_e)$.
    \item \textbf{Mass-polarization, diagonal:} $\sum_i\bar{\mathbf{p}}_i^{\,2}/(2M_N)$ --- modifies effective electron mass~\cite{bornoppenheimer1927,bunkerjensen}.
    \item \textbf{Mass-polarization, off-diagonal (two-electron):} $\sum_{i<l}\bar{\mathbf{p}}_i\!\cdot\!\bar{\mathbf{p}}_l/M_N$.
    \item \textbf{Coulomb:} $V_{NN}(Q)+V_{Ne}(Q,\{\bar{\mathbf{r}}_i\})+V_{ee}(\{\bar{\mathbf{r}}_i\})$. All Euler-angle independent.
\end{enumerate}

Compared to the previous (electrons-in-SF) framework: kinetic rot–electron coupling has \emph{replaced} potential rot–electron coupling. $V_{Ne}$ is now manifestly $\mathbf{S}$-independent --- the canonical chemistry advantage.\newpage

\section{Quantization}

The classical internal Hamiltonian (3.22) is parametrized by curvilinear coordinates (Euler angles $\phi,\theta,\chi$ and vibrational normal coordinates $Q_b$) and BF Cartesian electron coordinates $\bar{\mathbf{r}}_i$. Canonical quantization in curvilinear coordinates requires care: the naïve substitution $\pi\to-i\hbar\partial$ does not produce a Hermitian operator on the natural Hilbert space. The Podolsky (Laplace--Beltrami) prescription resolves this; the residual ordering produces the Watson pseudopotential, and the BF angular momentum acquires an \emph{anomalous} commutator algebra.

\subsection{The quantization problem in curvilinear coordinates}

Label the configuration-space coordinates $q^{A}$:
\begin{equation}
\{q^{A}\}\;=\;\bigl\{\,\phi,\theta,\chi,\;Q_1,\dots,Q_{3N_{\mathrm{nucl}}-6},\;\bar{\mathbf{r}}_1,\dots,\bar{\mathbf{r}}_{N_{\mathrm{elec}}}\,\bigr\}, \tag{4.1}
\end{equation}
with conjugate momenta $\pi_A$. The classical kinetic energy from (3.22) is a quadratic form
\begin{equation}
T \;=\; \tfrac12\,g^{AB}(q)\,\pi_A\pi_B, \tag{4.2}
\end{equation}
where $g^{AB}$ is the inverse of the configuration-space metric $g_{AB}$ defined by $T=\tfrac12 g_{AB}\dot q^A\dot q^B$. The metric inherits block structure from (3.22): a rotational block governed by the Watson reciprocal inertia $\boldsymbol{\mu}(Q)=(\mathbf{I}_N-\boldsymbol{\sigma}\boldsymbol{\sigma}^{\mathsf T})^{-1}$, an identity vibrational block (normal coordinates), and a Cartesian block for $\bar{\mathbf{r}}_i$.

\subsubsection{Podolsky form / Laplace--Beltrami quantization}

The natural Hilbert space is $L^{2}(\mathcal{C},\sqrt{|g|}\,d^{n}q)$, where the Riemannian volume element $\sqrt{|g|}\,d^{n}q$ is the unique measure compatible with the metric. To produce a Hermitian kinetic operator on this Hilbert space, replace $T$ by the Laplace--Beltrami operator divided by $-2/\hbar^{2}$:
\begin{equation}
\boxed{\;\hat T \;=\; -\frac{\hbar^{2}}{2}\,\frac{1}{\sqrt{|g|}}\,\partial_A\!\bigl(\sqrt{|g|}\,g^{AB}\,\partial_B\bigr).\;} \tag{4.3}
\end{equation}
This is the \emph{Podolsky prescription}~\cite{podolsky1928}. Equivalently, one may write $\hat\pi_A = -i\hbar(\partial_A + \tfrac12\partial_A\ln\sqrt{|g|})$ and then take $\hat T = \tfrac12\hat\pi_A g^{AB}\hat\pi_B$ in symmetric (Weyl-type) ordering; the two prescriptions coincide up to terms absorbed into the volume element.

For our problem, only the Euler-angle and $Q_b$ sectors are curvilinear; the BF Cartesian electronic sector retains the standard $\hat{\bar{\mathbf{p}}}_i = -i\hbar\nabla_{\bar{\mathbf{r}}_i}$ on $L^{2}(\mathbb{R}^{3N_{\mathrm{elec}}},d^{3N_{\mathrm{elec}}}\bar{\mathbf{r}})$.

\subsection{Volume element and Podolsky rescaling}

\subsubsection{Configuration-space metric and the volume element}

The configuration-space metric requires the kinetic energy with \emph{all} coordinates --- the electrons included --- carried as velocities $\mathbf{v}=(\boldsymbol{\omega},\dot{\mathbf{Q}},\dot{\bar{\mathbf{r}}}_i)$. (This velocity-form representation, used here solely to read off the measure, is distinct from the Routhian (3.13), in which the electrons are kept in momentum form.) Expanding the body-fixed nuclear, electronic, and mass-polarization kinetic terms via (3.9d)--(3.9e), in normal coordinates (vibrational block $\sum_j\mathbf{l}_{ja}\!\cdot\!\mathbf{l}_{jb}=\delta_{ab}$),
\begin{equation}
2T = \mathbf{v}^{\mathsf T}\mathsf{M}\,\mathbf{v},\qquad
\mathsf{M}=\begin{pmatrix}\mathbf{I} & \mathsf{Z} & \mathsf{C}\\[2pt] \mathsf{Z}^{\mathsf T} & \mathbb{1} & 0\\[2pt] \mathsf{C}^{\mathsf T} & 0 & \mathsf{D}\end{pmatrix},
\tag{4.3a}
\end{equation}
with $\mathbf{I}=\mathbf{I}_N+\mathbf{I}_e-\mathbf{I}_X$ the full instantaneous inertia ($\boldsymbol{\omega}$--$\boldsymbol{\omega}$ block), $\mathsf{Z}_{\alpha b}=(\boldsymbol{\sigma}_b)_\alpha$ the rotation--vibration Coriolis block, $\mathsf{C}$ the rotation--electron block (from $\boldsymbol{\omega}\!\cdot\!(\mathbf{L}_e-\mathbf{L}_X)$), and $\mathsf{D}_{(i\alpha)(j\beta)}=(m_e\delta_{ij}-m_e^2/M_{\mathrm{tot}})\,\delta_{\alpha\beta}$ the \emph{constant} electron block. The physical velocities relate to the coordinate velocities $\dot q=(\dot\phi,\dot\theta,\dot\chi,\dot Q_b,\dot{\bar{\mathbf{r}}}_i)$ by $\mathbf{v}=\mathsf{A}(q)\dot q$ with $\mathsf{A}=\mathrm{diag}\!\bigl(\mathsf{E}(\Omega),\mathbb{1},\mathbb{1}\bigr)$, where $\boldsymbol{\omega}=\mathsf{E}(\Omega)\dot{\boldsymbol{\Omega}}$ is the Euler kinematic relation. Hence $g_{AB}=(\mathsf{A}^{\mathsf T}\mathsf{M}\mathsf{A})_{AB}$ and
\begin{equation}
\det g = (\det\mathsf{A})^2\det\mathsf{M}=\sin^2\!\theta\;\det\mathsf{M},\qquad |\det\mathsf{E}(\Omega)|=\sin\theta . \tag{4.3b}
\end{equation}

Reduce $\det\mathsf{M}$ by the Schur complement of the constant electron block $\mathsf{D}$:
\begin{equation}
\det\mathsf{M}=\det\mathsf{D}\;\det\!\begin{pmatrix}\mathbf{I}-\mathsf{C}\mathsf{D}^{-1}\mathsf{C}^{\mathsf T} & \mathsf{Z}\\[2pt] \mathsf{Z}^{\mathsf T} & \mathbb{1}\end{pmatrix}. \tag{4.3c}
\end{equation}
The electron Schur correction is precisely the electronic + mass-polarization inertia $\mathbf{I}_e-\mathbf{I}_X$ of (3.9e): completing the square in $\dot{\bar{\mathbf{r}}}_i$ (equivalently, eliminating the electron velocities in favour of the momenta $\bar{\mathbf{p}}_i$) removes all $\boldsymbol{\omega}$-dependence from the electron block, $\mathsf{C}\mathsf{D}^{-1}\mathsf{C}^{\mathsf T}=\mathbf{I}_e-\mathbf{I}_X$, so $\mathbf{I}-\mathsf{C}\mathsf{D}^{-1}\mathsf{C}^{\mathsf T}=\mathbf{I}_N$ --- the same $\mathbf{I}_e/\mathbf{I}_X$ cancellation as the cross-check (3.18). A second Schur complement, now of the $\mathbb{1}$ vibrational block, gives
\begin{equation}
\det\mathsf{M}=\det\mathsf{D}\;\det\!\bigl(\mathbf{I}_N-\mathsf{Z}\mathsf{Z}^{\mathsf T}\bigr)=\det\mathsf{D}\;\det\mathbf{I}',\qquad \mathbf{I}'=\mathbf{I}_N-\boldsymbol{\sigma}\boldsymbol{\sigma}^{\mathsf T}, \tag{4.3d}
\end{equation}
using $\mathsf{Z}\mathsf{Z}^{\mathsf T}=\sum_b\boldsymbol{\sigma}_b\boldsymbol{\sigma}_b^{\mathsf T}=\boldsymbol{\sigma}\boldsymbol{\sigma}^{\mathsf T}$, with $\det\mathsf{D}$ a constant. The Coriolis coupling therefore enters the determinant through the \emph{same} modified inertia $\mathbf{I}'$ that defines the kernel $\boldsymbol{\mu}=(\mathbf{I}')^{-1}$. Dropping the constant electron factor (an overall normalization of the flat electronic measure), the Riemannian volume element is
\begin{equation}
\boxed{\;\sqrt{|g|}\,d^{n}q \;=\; \underbrace{\sin\theta\,d\phi\,d\theta\,d\chi}_{\text{Haar on }SO(3)}\cdot\underbrace{\sqrt{\det\mathbf{I}'(Q)}\,\prod_b dQ_b}_{\text{rot.\,--\,vib.\ Jacobian}}\cdot\underbrace{\prod_{i}d^{3}\bar{\mathbf{r}}_{i}}_{\text{flat (const.)}} . \;}\tag{4.4}
\end{equation}
The Jacobian carries the Watson modified inertia $\mathbf{I}'=\mathbf{I}_N-\boldsymbol{\sigma}\boldsymbol{\sigma}^{\mathsf T}$, \emph{not} the bare $\mathbf{I}_N$: the rotation--vibration Coriolis coupling reduces the rotational determinant from $\det\mathbf{I}_N$ to $\det(\mathbf{I}_N-\boldsymbol{\sigma}\boldsymbol{\sigma}^{\mathsf T})$.

\subsubsection{Wavefunction rescaling}

To work with the simpler measure $d\mu_W \equiv \sin\theta\,d\phi d\theta d\chi\prod_b dQ_b\prod_i d^{3}\bar{\mathbf{r}}_i$, define the rescaled wavefunction
\begin{equation}
\psi_W \;\equiv\; \bigl(\det\mathbf{I}'(Q)\bigr)^{1/4}\,\psi. \tag{4.5}
\end{equation}
Then $|\psi|^{2}\sqrt{|g|}\,d^n q = |\psi_W|^{2}\,d\mu_W$, so $\psi_W \in L^{2}(\mathcal{C},d\mu_W)$. The kinetic operator transforms by a similarity: $\hat T_W = \mathcal{J}^{1/2}\hat T \mathcal{J}^{-1/2}$ with $\mathcal{J} = (\det\mathbf{I}')^{1/2}$. The residual non-cancellation generates a new $Q$-dependent ordinary potential --- the Watson pseudopotential (next subsection).

\subsection{Anomalous BF angular-momentum algebra}

Total angular momentum is naturally defined as an SF observable; its BF components are obtained by projection through $\mathbf{S}$:
\begin{equation}
\hat J_{\alpha}^{\mathrm{BF}} \;=\; \mathbf{S}_{\beta\alpha}\,\hat J_\beta^{\mathrm{SF}}. \tag{4.6}
\end{equation}
The space-fixed components are ordinary angular-momentum operators in the Euler angles; they obey the normal algebra and act on the direction-cosine matrix $\mathbf{S}=\mathbf{S}(\Omega)$ through its \emph{space-fixed} index,
\begin{equation}
[\hat J_\gamma^{\mathrm{SF}},\hat J_\delta^{\mathrm{SF}}]=i\hbar\,\epsilon_{\gamma\delta\lambda}\hat J_\lambda^{\mathrm{SF}},\qquad
[\hat J_\gamma^{\mathrm{SF}},S_{\delta\beta}]=i\hbar\,\epsilon_{\gamma\delta\lambda}S_{\lambda\beta}. \tag{4.6a}
\end{equation}
Because $\mathbf{S}$ is a proper rotation ($\det\mathbf{S}=1$), the triple product of direction cosines collapses to one Levi-Civita symbol, and orthogonality inverts the projection (4.6):
\begin{equation}
\epsilon_{\gamma\delta\lambda}\,S_{\gamma\alpha}S_{\delta\beta}S_{\lambda\nu}=\det(\mathbf{S})\,\epsilon_{\alpha\beta\nu}=\epsilon_{\alpha\beta\nu},\qquad
\hat J_\lambda^{\mathrm{SF}}=S_{\lambda\nu}\,\hat J_\nu^{\mathrm{BF}}. \tag{4.6b}
\end{equation}
Now evaluate the body-fixed commutator (writing $\hat J\equiv\hat J^{\mathrm{SF}}$), moving the Euler-angle functions $S_{\delta\beta}$ leftward past $\hat J_\gamma$ with (4.6a):
\begin{align*}
[\hat J_\alpha^{\mathrm{BF}},\hat J_\beta^{\mathrm{BF}}]
&=[S_{\gamma\alpha}\hat J_\gamma,\,S_{\delta\beta}\hat J_\delta]
= S_{\gamma\alpha}S_{\delta\beta}[\hat J_\gamma,\hat J_\delta]
+ S_{\gamma\alpha}\,[\hat J_\gamma,S_{\delta\beta}]\,\hat J_\delta
- S_{\delta\beta}\,[\hat J_\delta,S_{\gamma\alpha}]\,\hat J_\gamma\\
&= i\hbar\,\epsilon_{\gamma\delta\lambda}\bigl(S_{\gamma\alpha}S_{\delta\beta}\hat J_\lambda+S_{\gamma\alpha}S_{\lambda\beta}\hat J_\delta+S_{\delta\beta}S_{\lambda\alpha}\hat J_\gamma\bigr),
\end{align*}
where $[\hat J_\gamma,\hat J_\delta]=i\hbar\epsilon_{\gamma\delta\lambda}\hat J_\lambda$, $[\hat J_\gamma,S_{\delta\beta}]=i\hbar\epsilon_{\gamma\delta\lambda}S_{\lambda\beta}$, and $-[\hat J_\delta,S_{\gamma\alpha}]=-i\hbar\epsilon_{\delta\gamma\lambda}S_{\lambda\alpha}=+i\hbar\epsilon_{\gamma\delta\lambda}S_{\lambda\alpha}$ were collected under a common $\epsilon_{\gamma\delta\lambda}$. Insert $\hat J=\mathbf{S}\hat J^{\mathrm{BF}}$ in each term and apply (4.6b):
\begin{equation*}
\epsilon_{\gamma\delta\lambda}S_{\gamma\alpha}S_{\delta\beta}\,\hat J_\lambda=\epsilon_{\alpha\beta\nu}\hat J_\nu^{\mathrm{BF}},\quad
\epsilon_{\gamma\delta\lambda}S_{\gamma\alpha}S_{\lambda\beta}\,\hat J_\delta=\epsilon_{\alpha\nu\beta}\hat J_\nu^{\mathrm{BF}},\quad
\epsilon_{\gamma\delta\lambda}S_{\delta\beta}S_{\lambda\alpha}\,\hat J_\gamma=\epsilon_{\nu\beta\alpha}\hat J_\nu^{\mathrm{BF}}.
\end{equation*}
Since $\epsilon_{\alpha\beta\nu}+\epsilon_{\alpha\nu\beta}+\epsilon_{\nu\beta\alpha}=\epsilon_{\alpha\beta\nu}-\epsilon_{\alpha\beta\nu}-\epsilon_{\alpha\beta\nu}=-\epsilon_{\alpha\beta\nu}$, the three pieces sum to a single term of \emph{reversed} sign --- the extra contribution from differentiating $\mathbf{S}$ outweighs and flips the bare SF algebra:
\begin{equation}
\boxed{\;[\hat J_{\alpha}^{\mathrm{BF}},\hat J_{\beta}^{\mathrm{BF}}]\;=\;-i\hbar\,\epsilon_{\alpha\beta\gamma}\,\hat J_{\gamma}^{\mathrm{BF}}.\;} \tag{4.7}
\end{equation}
This is the famous \textbf{anomalous commutator} for BF angular momentum~\cite{klein1929,vanvleck1951}, reviewed in Bunker--Jensen~\cite{bunkerjensen}.

In contrast, BF electronic angular momentum is built intrinsically from BF Cartesian electron operators:
\begin{equation}
[\hat L_{\mathrm{elec},\alpha},\hat L_{\mathrm{elec},\beta}]\;=\;+i\hbar\,\epsilon_{\alpha\beta\gamma}\,\hat L_{\mathrm{elec},\gamma}, \tag{4.8}
\end{equation}
the \emph{normal} algebra. Since $\hat J^{\mathrm{BF}}$ acts only on Euler angles and $\hat{\mathbf{L}}_{\mathrm{elec}}$ acts only on BF electron coordinates, the two commute:
\begin{equation}
[\hat J_{\alpha}^{\mathrm{BF}},\hat L_{\mathrm{elec},\beta}] \;=\; 0. \tag{4.9}
\end{equation}
The vibrational $\hat\pi_\alpha$, built from $Q_b$ and $\hat P_b$, also commutes with both $\hat J^{\mathrm{BF}}$ and $\hat{\mathbf{L}}_{\mathrm{elec}}$. From here on we write $\hat J_\alpha \equiv \hat J_\alpha^{\mathrm{BF}}$ and the anomalous algebra (4.7) is understood.

\subsection{Watson pseudopotential}

We now quantize the rovibrational kinetic energy of (3.22),
\begin{equation}
T_{\mathrm{rovib}}=\tfrac12\sum_{\alpha\beta}A_\alpha\,\mu_{\alpha\beta}(Q)\,A_\beta+\tfrac12\sum_b P_b^2,\qquad
A_\alpha\equiv J_\alpha-L_{\mathrm{elec},\alpha}-\pi_\alpha, \tag{4.10a}
\end{equation}
a quadratic form in the momenta $(J_\alpha,P_b)$ whose coefficients $\mu_{\alpha\beta}(Q)$ depend on the shape. The na\"ive replacement $P_b\to-i\hbar\,\partial_{Q_b}$ is \emph{not} Hermitian on the curved measure (4.4); the Podolsky (Laplace--Beltrami) prescription (4.3) restores Hermiticity, and the rescaling (4.5) to the flat measure then leaves a $c$-number residue --- the \textbf{Watson pseudopotential} $\hat U(Q)$. We carry out the reduction in full.

\subsubsection{The curvilinear extra-potential, in general}

For arbitrary coordinates $q^A$ with covariant metric $g_{AB}$ ($2T=g_{AB}\dot q^A\dot q^B$), contravariant inverse $g^{AB}$, and determinant $g\equiv\det g_{AB}$, the kinetic operator Hermitian on $L^2(\sqrt g\,d^nq)$ is the Laplace--Beltrami operator
\begin{equation}
\hat T_{\mathrm{LB}}=-\frac{\hbar^2}{2}\,g^{-1/2}\,\partial_A\!\bigl(g^{1/2}\,g^{AB}\,\partial_B\bigr). \tag{4.10b}
\end{equation}
To pass to the flat measure $\prod_A dq^A$, set $\psi=g^{-1/4}\psi_W$ (so $\psi_W=g^{1/4}\psi$) and conjugate, $\hat T_W=g^{1/4}\hat T_{\mathrm{LB}}\,g^{-1/4}$. Writing $\rho\equiv g^{1/2}$ one has $\partial_B\psi=\rho^{-1/2}f_B$ with $f_B\equiv\partial_B\psi_W-\tfrac12(\partial_B\ln\rho)\psi_W$; carrying the $g^{1/4}$ through $\hat T_{\mathrm{LB}}$,
\begin{equation}
\hat T_W\psi_W=-\frac{\hbar^2}{2}\Bigl[\underbrace{\tfrac12(\partial_A\ln\rho)\,g^{AB}f_B}_{(\mathrm i)}+\underbrace{(\partial_Ag^{AB})\,f_B}_{(\mathrm{ii})}+\underbrace{g^{AB}\,\partial_Af_B}_{(\mathrm{iii})}\Bigr]. \tag{4.10c}
\end{equation}
Insert $f_B$ and $\partial_Af_B=\partial_A\partial_B\psi_W-\tfrac12(\partial_A\partial_B\ln\rho)\psi_W-\tfrac12(\partial_B\ln\rho)\partial_A\psi_W$, and sort by the order of derivative acting on $\psi_W$:
\begin{itemize}
\item[$\partial\partial\psi_W:$] only (iii) contributes, $g^{AB}\partial_A\partial_B\psi_W$;
\item[$\partial\psi_W:$] (i) gives $+\tfrac12 g^{AB}(\partial_A\ln\rho)\partial_B\psi_W$ and (iii) gives $-\tfrac12 g^{AB}(\partial_B\ln\rho)\partial_A\psi_W$, which \emph{cancel} by symmetry of $g^{AB}$, leaving only $(\partial_Ag^{AB})\partial_B\psi_W$ from (ii);
\item[$\psi_W:$] (i) gives $-\tfrac14 g^{AB}(\partial_A\ln\rho)(\partial_B\ln\rho)$, (ii) gives $-\tfrac12(\partial_Ag^{AB})(\partial_B\ln\rho)$, (iii) gives $-\tfrac12 g^{AB}\partial_A\partial_B\ln\rho$.
\end{itemize}
The first two derivative-orders reassemble into $\partial_A(g^{AB}\partial_B\psi_W)$, i.e.\ the flat symmetric operator $\tfrac12\hat p_Ag^{AB}\hat p_B$ ($\hat p_A=-i\hbar\partial_A$). The $\psi_W$-terms, using $-\tfrac12(\partial_Ag^{AB})(\partial_B\ln\rho)-\tfrac12 g^{AB}\partial_A\partial_B\ln\rho=-\tfrac12\partial_A(g^{AB}\partial_B\ln\rho)$ and $\ln\rho=\tfrac12\ln g$, give the extra-potential:
\begin{equation}
\boxed{\;\hat T_W=\tfrac12\,\hat p_A\,g^{AB}\,\hat p_B+U,\qquad
U=\frac{\hbar^2}{8}\Bigl[\tfrac14\,g^{AB}(\partial_A\ln g)(\partial_B\ln g)+\partial_A\!\bigl(g^{AB}\partial_B\ln g\bigr)\Bigr].\;}\tag{4.10d}
\end{equation}
Equation (4.10d) is exact and coordinate-general; the prefactor $\hbar^2/8$ traces to the two half-powers $g^{\pm1/4}$ of the rescaling.

\subsubsection{The molecular metric and its determinant}

Read the contravariant metric off (4.10a). With $\pi_\alpha=\sum_b\sigma_{b\alpha}P_b$ (i.e.\ $\boldsymbol\pi=\boldsymbol\sigma^{\mathsf T}\!P$, $\sigma_{b\alpha}=(\boldsymbol\sigma_b)_\alpha$), and dropping the electronic shift $\hat{\mathbf L}_{\mathrm{elec}}$ (a different sector, commuting with the nuclear metric),
\begin{equation*}
2T_{\mathrm{rovib}}=J^{\mathsf T}\boldsymbol\mu J-2J^{\mathsf T}\boldsymbol\mu\boldsymbol\sigma^{\mathsf T}P+P^{\mathsf T}(\mathbb 1+\boldsymbol\sigma\boldsymbol\mu\boldsymbol\sigma^{\mathsf T})P,
\end{equation*}
so in the body-frame momentum basis $(J_\alpha,P_b)$,
\begin{equation}
g^{AB}=\begin{pmatrix}\boldsymbol\mu & -\boldsymbol\mu\boldsymbol\sigma^{\mathsf T}\\[2pt] -\boldsymbol\sigma\boldsymbol\mu & \mathbb 1+\boldsymbol\sigma\boldsymbol\mu\boldsymbol\sigma^{\mathsf T}\end{pmatrix}. \tag{4.10e}
\end{equation}
By the Schur complement of the $\boldsymbol\mu$ block,
\begin{equation}
\det g^{AB}=\det\boldsymbol\mu\;\det\!\bigl(\mathbb 1+\boldsymbol\sigma\boldsymbol\mu\boldsymbol\sigma^{\mathsf T}-\boldsymbol\sigma\boldsymbol\mu\,\boldsymbol\mu^{-1}\boldsymbol\mu\boldsymbol\sigma^{\mathsf T}\bigr)=\det\boldsymbol\mu, \tag{4.10f}
\end{equation}
hence $g=\det g_{AB}=1/\det\boldsymbol\mu=\det\mathbf I'\equiv\gamma(Q)$, consistent with $\sqrt{|g|}=\sin\theta\sqrt{\gamma}$ of (4.4). The full measure carries $\ln g=2\ln\sin\theta+\ln\gamma$. In the rovibrational convention the body-frame $\hat J_\alpha$ are already Hermitian on $\sin\theta\,d\Omega$ [with the anomalous algebra (4.7)], so the $2\ln\sin\theta$ is accounted for by the $\hat J_\alpha$; \emph{only} $\gamma(Q)$ is removed by the rescaling (4.5), and we set $\ln g\to\ln\gamma$ (a function of $Q$ alone) in (4.10d).

\subsubsection{Rotational sector: no $c$-number}

The only operator ambiguity in the rotational kernel is the ordering of $\hat J_\alpha\hat J_\beta$. Splitting into symmetric and antisymmetric parts and using the anomalous algebra (4.7),
\begin{equation}
\tfrac12\mu_{\alpha\beta}\hat J_\alpha\hat J_\beta=\tfrac14\mu_{\alpha\beta}\{\hat J_\alpha,\hat J_\beta\}+\tfrac14\mu_{\alpha\beta}[\hat J_\alpha,\hat J_\beta]=\tfrac14\mu_{\alpha\beta}\{\hat J_\alpha,\hat J_\beta\}-\frac{i\hbar}{4}\,\mu_{\alpha\beta}\epsilon_{\alpha\beta\gamma}\hat J_\gamma. \tag{4.10g}
\end{equation}
Since $\mu_{\alpha\beta}=\mu_{\beta\alpha}$ is symmetric while $\epsilon_{\alpha\beta\gamma}$ is antisymmetric in $(\alpha,\beta)$, the contraction $\mu_{\alpha\beta}\epsilon_{\alpha\beta\gamma}=0$: the commutator term vanishes, the asymmetric-top operator is Hermitian as written, and it contributes \emph{no} $c$-number. The same holds for the full kernel $\tfrac12\hat A_\alpha\mu_{\alpha\beta}\hat A_\beta$ ($\hat A_\alpha=\hat J_\alpha-\hat L_{\mathrm{elec},\alpha}-\hat\pi_\alpha$), since $\boldsymbol\mu$ is symmetric and commutes with $\hat J$ and $\hat{\mathbf L}_{\mathrm{elec}}$.

\subsubsection{Vibrational sector: the rescaling, step by step}

With $g^{ab}=\delta_{ab}$ but the measure carrying $\gamma(Q)$, the vibrational part of (4.10b) is
\begin{equation}
\hat T_{\mathrm{vib}}=-\frac{\hbar^2}{2}\,\gamma^{-1/2}\sum_b\partial_{Q_b}\!\bigl(\gamma^{1/2}\,\partial_{Q_b}\bigr). \tag{4.10h}
\end{equation}
Conjugate by $\gamma^{1/4}$ and act on $\psi_W$, working from the innermost derivative outward (each line is one application of the product rule):
\begin{align}
\partial_{Q_b}\!\bigl(\gamma^{-1/4}\psi_W\bigr)&=\gamma^{-1/4}\Bigl[\partial_{Q_b}\psi_W-\tfrac14(\partial_{Q_b}\!\ln\gamma)\psi_W\Bigr],\nonumber\\
\gamma^{1/2}\,\partial_{Q_b}\!\bigl(\gamma^{-1/4}\psi_W\bigr)&=\gamma^{1/4}\Bigl[\partial_{Q_b}\psi_W-\tfrac14(\partial_{Q_b}\!\ln\gamma)\psi_W\Bigr],\nonumber\\
\partial_{Q_b}\Bigl(\gamma^{1/4}\bigl[\,\cdots\,\bigr]\Bigr)&=\gamma^{1/4}\Bigl[\partial_{Q_b}^2\psi_W-\tfrac1{16}(\partial_{Q_b}\!\ln\gamma)^2\psi_W-\tfrac14(\partial_{Q_b}^2\ln\gamma)\psi_W\Bigr], \tag{4.10i}
\end{align}
where in the last line the two cross terms $+\tfrac14(\partial_{Q_b}\!\ln\gamma)\partial_{Q_b}\psi_W$ (from differentiating $\gamma^{1/4}$) and $-\tfrac14(\partial_{Q_b}\!\ln\gamma)\partial_{Q_b}\psi_W$ (from the bracket) cancel. Multiplying by $-\tfrac{\hbar^2}{2}\,\gamma^{1/4}\!\cdot\!\gamma^{-1/2}=-\tfrac{\hbar^2}{2}\gamma^{-1/4}$ and summing over $b$,
\begin{equation}
\gamma^{1/4}\hat T_{\mathrm{vib}}\gamma^{-1/4}=\tfrac12\sum_b\hat P_b^2+U_{\mathrm{vib}},\qquad \hat P_b=-i\hbar\,\partial_{Q_b}, \tag{4.10j}
\end{equation}
with the \emph{Pauli residue}
\begin{equation}
U_{\mathrm{vib}}=\frac{\hbar^2}{8}\sum_b\Bigl[\partial_{Q_b}^2\ln\gamma+\tfrac14(\partial_{Q_b}\!\ln\gamma)^2\Bigr]=\frac{\hbar^2}{2}\,\gamma^{-1/4}\Bigl(\sum_b\partial_{Q_b}^2\Bigr)\gamma^{1/4}. \tag{4.10k}
\end{equation}
This is exactly what (4.10d) yields for the \emph{flat} vibrational block $g^{ab}=\delta_{ab}$; the off-diagonal $-\boldsymbol\mu\boldsymbol\sigma^{\mathsf T}$ and the $\boldsymbol\sigma\boldsymbol\mu\boldsymbol\sigma^{\mathsf T}$ dressing of (4.10e) supply the remaining \emph{Coriolis} $c$-numbers. Collecting the three sectors,
\begin{align}
\hat T_{\mathrm{rovib}}\;=\;&\tfrac12\!\sum_{\alpha\beta}\bigl(\hat J_\alpha-\hat L_{\mathrm{elec},\alpha}-\hat\pi_\alpha\bigr)\mu_{\alpha\beta}(Q)\bigl(\hat J_\beta-\hat L_{\mathrm{elec},\beta}-\hat\pi_\beta\bigr)\nonumber\\
&+\tfrac12\!\sum_b\hat P_b^2\;+\;\hat U(Q),\qquad \hat U=U_{\mathrm{vib}}+U_{\mathrm{Cor}}. \tag{4.10}
\end{align}

\subsubsection{Collapse to the Watson trace}

The residue (4.10k) and the Coriolis $c$-numbers $U_{\mathrm{Cor}}$ collapse by two exact identities. First, Jacobi's formula for $\gamma=\det\mathbf I'$ and the derivative of an inverse,
\begin{equation}
\partial_b\ln\gamma=\operatorname{tr}\!\bigl(\boldsymbol\mu\,\partial_b\mathbf I'\bigr)=\sum_{\alpha\beta}\mu_{\alpha\beta}\,\partial_b\mathbf I'_{\alpha\beta},\qquad \partial_b\boldsymbol\mu=-\,\boldsymbol\mu\,(\partial_b\mathbf I')\,\boldsymbol\mu, \tag{4.10l}
\end{equation}
re-express every $\partial_b\ln\gamma$ and $\partial_b\boldsymbol\mu$ through $\partial_b\mathbf I'$. Second, since $\mathbf I'=\mathbf I_N-\boldsymbol\sigma\boldsymbol\sigma^{\mathsf T}$ with $\boldsymbol\sigma_b=\sum_a\boldsymbol\zeta_{ab}Q_a$ \emph{linear} in $Q$ [(3.17a)],
\begin{equation}
\partial_c\mathbf I'_{\alpha\beta}=-\sum_b\bigl(\zeta^\alpha_{cb}\,\sigma_{b\beta}+\sigma_{b\alpha}\,\zeta^\beta_{cb}\bigr),\qquad \zeta^\alpha_{cb}\equiv(\boldsymbol\zeta_{cb})_\alpha=\partial_c\sigma_{b\alpha}\ (\text{constant}), \tag{4.10m}
\end{equation}
so $\partial_b\mathbf I'$ is built from the \emph{constant} Coriolis matrices $\boldsymbol\zeta_{cb}=\sum_j\mathbf l_{jc}\!\times\!\mathbf l_{jb}$. Substituting (4.10l)--(4.10m) into the vibrational residue (4.10k) and into $U_{\mathrm{Cor}}$, the gradient terms $\propto\partial_b\operatorname{tr}\boldsymbol\mu$ generated by the two groups cancel identically; the surviving $Q$-local remainder is fixed by the Coriolis sum rules obeyed by the $\boldsymbol\zeta$'s [from the orthonormality $\sum_j\mathbf l_{ja}\!\cdot\!\mathbf l_{jb}=\delta_{ab}$ of the Eckart mode vectors, the $\boldsymbol\zeta^{(\alpha)}$ acting as $SO(3)$ generators on the mode index], which collapse it to a single trace over the three Cartesian components~\cite{watson1968}:
\begin{equation}
\boxed{\;\hat U(Q)\;=\;-\frac{\hbar^{2}}{8}\sum_{\alpha=1}^{3}\mu_{\alpha\alpha}(Q),\qquad \boldsymbol{\mu}=(\mathbf{I}_N-\boldsymbol{\sigma}\boldsymbol{\sigma}^{\mathsf T})^{-1}.\;} \tag{4.11}
\end{equation}
Two remarks. \emph{(a)} The cancellation is special to \emph{rectilinear} normal coordinates: for general curvilinear $Q$ the linear relation (4.10m) fails, the residue (4.10k) does \emph{not} collapse to the trace, and additional vibrational extrapotential terms survive. \emph{(b)} The electrons do not modify $\hat U$: $\mu_{\alpha\beta}$ depends on $Q$ alone, the electronic sector enters the rovib block only through $\hat{\mathbf{L}}_{\mathrm{elec}}$ (which commutes with $\boldsymbol\mu$ and with the $Q$-measure $\gamma$), and the electronic Cartesian Jacobian is flat. The pseudopotential is a $c$-number on the nuclear configuration space, fixed entirely by the modified inertia $\mathbf{I}'$.

\subsubsection{Algebraic cross-check of the coefficient}

The $-\hbar^2/8$ is confirmed by the equivalent algebraic (Weyl-ordering) route noted under (4.3): with the Hermitian momenta $\hat\pi_A=-i\hbar(\partial_A+\tfrac12\partial_A\ln\sqrt{|g|})$ one has $\hat T=\tfrac12\hat\pi_A g^{AB}\hat\pi_B$, identical to the Laplace--Beltrami operator. Three observations fix the result:
\begin{itemize}
\item \emph{Vibrational sector.} With $g^{ab}=\delta_{ab}$ and $\sqrt{|g|}\propto\gamma^{1/2}$, expanding $\tfrac12\sum_b\hat\pi_{Q_b}^2$ (the Hermitian momenta of (4.3) specialized to the modes) and rescaling reproduces exactly the residue (4.10k), $U_{\mathrm{vib}}=\tfrac{\hbar^2}{2}\gamma^{-1/4}\bigl(\sum_b\partial_{Q_b}^2\bigr)\gamma^{1/4}$ --- the $-\hbar^2/8$ scale enters through the two $\tfrac12\partial_A\ln\sqrt{|g|}$ factors, $\tfrac12\hat\pi_{Q_b}^2\supset\tfrac{\hbar^2}{8}(\partial_{Q_b}\ln\gamma)^2$ etc.
\item \emph{Rotational sector.} $\mu_{\alpha\beta}(Q)$ commutes with $\hat J$, so the only reordering is among the $\hat J$'s; by the anomalous algebra (4.7) the antisymmetric part of $\hat J_\alpha\mu_{\alpha\beta}\hat J_\beta$ vanishes ($\mu_{\alpha\beta}=\mu_{\beta\alpha}$), so this block adds no independent $c$-number.
\item \emph{Reduction.} Carrying the Watson cancellations (Jacobi's formula and the Coriolis sum rules) through both sectors collapses every $Q$-derivative to the single trace $\hat U=-\tfrac{\hbar^2}{8}\sum_\alpha\mu_{\alpha\alpha}$, in agreement with (4.11).
\end{itemize}
The pseudopotential is a $c$-number depending only on $Q$ (a function on the nuclear configuration space).

\subsubsection{Operator ordering: what is and isn't ambiguous}

\begin{itemize}
\item $\mu_{\alpha\beta}(Q)$ commutes with $\hat J$ (different coordinate sectors: Euler angles vs.\ $Q$), with $\hat{\mathbf{L}}_{\mathrm{elec}}$ (BF electron coords), but not with $\hat P_b$ (since $\partial\mu_{\alpha\beta}/\partial Q_b \ne 0$). Hence $\hat J_\alpha\mu_{\alpha\beta}\hat J_\beta = \mu_{\alpha\beta}\hat J_\alpha\hat J_\beta$.
\item Within the rotational block, the symmetric product $\hat J_\alpha\mu_{\alpha\beta}\hat J_\beta$ in (4.10) is unambiguous because $\mu$ is symmetric in $(\alpha,\beta)$, and the Watson pseudopotential (4.11) accounts for the residual non-commutativity of the $\hat J$'s themselves.
\item In normal coordinates the vibrational term is the bare $\tfrac12\sum_b\hat P_b^2$; the $Q$-dependence of the volume element has been moved into $\hat U(Q)$ by the rescaling (4.5), which is precisely what Podolsky produces.
\end{itemize}

\subsection{Electronic kinetic operator}

In BF Cartesian electron coordinates, $\hat{\bar{\mathbf{p}}}_i = -i\hbar\nabla_{\bar{\mathbf{r}}_i}$ is Hermitian on $L^{2}(\mathbb{R}^{3N_{\mathrm{elec}}},d^{3N_{\mathrm{elec}}}\bar{\mathbf{r}})$. The electronic + mass-polarization kinetic operator is therefore the straightforward quantization of (3.21):
\begin{equation}
\hat T_e \;=\; \sum_i\frac{\hat{\bar{\mathbf{p}}}_i^{\,2}}{2m_e} \;+\; \frac{1}{2M_N}\Bigl(\sum_i\hat{\bar{\mathbf{p}}}_i\Bigr)^{\!2}. \tag{4.12}
\end{equation}
Expanding the second piece:
\begin{equation}
\frac{1}{2M_N}\Bigl(\sum_i\hat{\bar{\mathbf{p}}}_i\Bigr)^{\!2} \;=\; \sum_i\frac{\hat{\bar{\mathbf{p}}}_i^{\,2}}{2M_N} \;+\; \frac{1}{M_N}\!\sum_{i<l}\hat{\bar{\mathbf{p}}}_i\!\cdot\!\hat{\bar{\mathbf{p}}}_l. \tag{4.13}
\end{equation}
The diagonal $1/(2M_N)$ piece renormalizes the effective electron mass through $1/m_e + 1/M_N = M_{\mathrm{tot}}/(m_e M_N) = 1/\mu_e$. The off-diagonal two-electron term is the mass-polarization \emph{sensu stricto}, dynamically coupling electrons through nuclear recoil --- a purely classical effect (Section 2.4) inherited by the quantum theory.

\subsection{Complete quantum molecular Hamiltonian}

Combining (4.10), (4.11), (4.12), and the rotationally invariant Coulomb potentials from (3.8):
\begin{equation}
\boxed{
\begin{aligned}
\hat H_{\mathrm{int}}\;=\;& \tfrac12\sum_{\alpha\beta}\bigl(\hat J_\alpha - \hat L_{\mathrm{elec},\alpha} - \hat\pi_\alpha\bigr)\mu_{\alpha\beta}(Q)\bigl(\hat J_\beta - \hat L_{\mathrm{elec},\beta} - \hat\pi_\beta\bigr) \\
& + \tfrac12\sum_b\hat P_b^2 \\
& + \sum_i\frac{\hat{\bar{\mathbf{p}}}_i^{\,2}}{2m_e} \;+\; \frac{1}{2M_N}\Bigl(\sum_i\hat{\bar{\mathbf{p}}}_i\Bigr)^{\!2} \\
& - \frac{\hbar^{2}}{8}\sum_{\alpha}\mu_{\alpha\alpha}(Q) \\
& + V_{NN}(Q) \;+\; V_{Ne}(Q,\{\bar{\mathbf{r}}_i\}) \;+\; V_{ee}(\{\bar{\mathbf{r}}_i\}),
\end{aligned}}\tag{4.14}
\end{equation}
acting on $L^{2}(\mathcal{C},\,d\mu_W)$ with $d\mu_W = \sin\theta\,d\phi d\theta d\chi\,\prod_b dQ_b\,\prod_i d^{3}\bar{\mathbf{r}}_i$. All operators are Hermitian on this Hilbert space.

This is the \textbf{exact non-relativistic molecular Hamiltonian} after separation of the total COM motion, in BF + Eckart frame, with electrons referred to the nuclear COM. The only approximations made along the way were: (i) non-relativistic dynamics; (ii) Eckart's choice of the BF orientation (a convention, not an approximation); (iii) parametrization of the nuclear shape by $3N_{\mathrm{nucl}}-6$ internal coordinates $Q_b$ (a parametrization, also exact).

\subsection{Born--Huang expansion}

\subsubsection{Clamped-nuclei electronic Hamiltonian}

At each nuclear configuration $Q$, define the \textbf{BF clamped-nuclei electronic Hamiltonian} as the sum of all $Q$-parametric electron-only pieces of (4.14):
\begin{equation}
\hat H_{\mathrm{el}}(Q) \;\equiv\; \sum_i\frac{\hat{\bar{\mathbf{p}}}_i^{\,2}}{2m_e} \;+\; \frac{1}{2M_N}\Bigl(\sum_i\hat{\bar{\mathbf{p}}}_i\Bigr)^{\!2} \;+\; V_{Ne}(Q,\{\bar{\mathbf{r}}_i\}) \;+\; V_{ee}(\{\bar{\mathbf{r}}_i\}). \tag{4.15}
\end{equation}
\emph{Remark.} The mass-polarization piece $(\sum_i\hat{\bar{\mathbf{p}}}_i)^{2}/(2M_N)$ is included here as a small correction with $1/M_N$. Some treatments (``clamped-nuclei BO'') drop it; we keep it explicit because it is parametric in $Q$ only through the absent dependence (i.e., $M_N$ is just a number), which means it acts trivially on the parametric family of electronic Hilbert spaces.

Solve the parametric eigenvalue problem:
\begin{equation}
\hat H_{\mathrm{el}}(Q)\,|\phi_n(Q)\rangle \;=\; E_n^{\mathrm{el}}(Q)\,|\phi_n(Q)\rangle, \tag{4.16}
\end{equation}
in the electronic Hilbert space $\mathcal{H}_e = L^{2}(\mathbb{R}^{3N_{\mathrm{elec}}},d^{3N_{\mathrm{elec}}}\bar{\mathbf{r}})$. The eigenstates $\{|\phi_n(Q)\rangle\}$ depend \emph{parametrically} on $Q$ --- $Q$ is treated as a fixed classical parameter, not a quantum degree of freedom --- and form a complete orthonormal basis of $\mathcal{H}_e$ at each $Q$:
\begin{equation}
\sum_n |\phi_n(Q)\rangle\langle\phi_n(Q)| \;=\; \hat{\mathbb{1}}_{e},\qquad \langle\phi_m(Q)|\phi_n(Q)\rangle_e = \delta_{mn}. \tag{4.17}
\end{equation}

\subsubsection{The exact Born--Huang expansion}

The \emph{exact} molecular wavefunction $|\Psi\rangle \in L^{2}(\mathcal{C},d\mu_W)$ admits a resolution against the parametric electronic basis. Insert the resolution of identity (4.17) at each $Q$:
\begin{equation}
|\Psi\rangle \;=\; \sum_n |\chi_n\rangle \otimes |\phi_n(Q)\rangle, \qquad \chi_n(\Omega,Q) \;\equiv\; \langle\phi_n(Q)|\Psi\rangle_e, \tag{4.18}
\end{equation}
where $\Omega=(\phi,\theta,\chi)$, the electronic inner product is taken at fixed $Q$, and $\chi_n$ are the \textbf{nuclear coefficient functions}. The expansion (4.18) is \textbf{exact}, not an ansatz: it follows from the completeness (4.17) at each $Q$, with the only subtlety being that the basis $|\phi_n(Q)\rangle$ varies parametrically. There is \emph{no} simple tensor factorization $\mathcal{H}_{\mathrm{molecule}} = \mathcal{H}_N \otimes \mathcal{H}_e$ because of this parametric dependence; the Born--Huang expansion~\cite{bornhuang1954} is the rigorous replacement.

\subsubsection{Coupled-channel equations}

Substitute (4.18) into $\hat H_{\mathrm{int}}|\Psi\rangle = E|\Psi\rangle$, project from the left by $\langle\phi_m(Q)|$, and use (4.16):
\begin{equation}
\bigl[\hat T_N + E_n^{\mathrm{el}}(Q) + V_{NN}(Q) + \hat U(Q)\bigr]\chi_n(\Omega,Q) \;+\; \sum_m \hat\Lambda_{nm}\,\chi_m(\Omega,Q) \;=\; E\,\chi_n(\Omega,Q), \tag{4.19}
\end{equation}
where $\hat T_N$ denotes the nuclear (rotational + vibrational) kinetic operator from the first two lines of (4.14). The \textbf{non-adiabatic coupling operators} $\hat\Lambda_{nm}$ arise from the action of $\hat T_N$ (and the rot--electronic couplings through $\hat{\mathbf{L}}_{\mathrm{elec}}$) on the $Q$-parametric dependence of $|\phi_m(Q)\rangle$:
\begin{equation}
\hat\Lambda_{nm} \;\sim\; -\frac{\hbar^{2}}{1}\,\langle\phi_n|\nabla_Q\phi_m\rangle_e\!\cdot\!\nabla_Q \;-\; \frac{\hbar^{2}}{2}\,\langle\phi_n|\nabla_Q^{2}\phi_m\rangle_e \;+\;\ldots, \tag{4.20}
\end{equation}
with analogous terms from the rotation--electron Coriolis $\hat J_\alpha\hat L_{\mathrm{elec},\alpha}$ kernel.

\subsubsection{Born--Oppenheimer approximation as truncation}

The \textbf{Born--Oppenheimer (BO) approximation} is the diagonal truncation of (4.19): set $\hat\Lambda_{nm}\to 0$ for $m\ne n$ and (optionally) also $\hat\Lambda_{nn}\to 0$. The resulting effective nuclear problem is
\begin{equation}
\bigl[\hat T_N + E_n^{\mathrm{el}}(Q) + V_{NN}(Q) + \hat U(Q)\bigr]\chi_n^{\mathrm{BO}}(\Omega,Q) \;=\; E^{\mathrm{BO}}\,\chi_n^{\mathrm{BO}}(\Omega,Q), \tag{4.21}
\end{equation}
where the \textbf{Born--Oppenheimer potential energy surface} for electronic state $n$ is
\begin{equation}
V_n^{\mathrm{BO}}(Q) \;\equiv\; E_n^{\mathrm{el}}(Q) + V_{NN}(Q) + \hat U(Q). \tag{4.22}
\end{equation}
The exactness of (4.18) and the well-definedness of (4.19) are mathematical facts. BO is the physical assumption that off-diagonal $\hat\Lambda_{nm}$ are small --- justified when electronic gaps are large compared to nuclear kinetic energies, breaking down at conical intersections and avoided crossings, where the parametric variation of $|\phi_n(Q)\rangle$ becomes singular.

\subsubsection{Summary}

\begin{enumerate}
\item (4.14) is the exact non-relativistic quantum molecular Hamiltonian in BF + Eckart + nuclear-COM coordinates.
\item The Born--Huang expansion (4.18) is exact, following from completeness of the parametric electronic basis (4.17).
\item The infinite coupled system (4.19) is equivalent to the Schr\"odinger equation $\hat H_{\mathrm{int}}|\Psi\rangle = E|\Psi\rangle$.
\item The Born--Oppenheimer truncation reduces it to (4.21), giving the familiar single-PES picture (4.22).
\end{enumerate}


\newpage

\section*{Appendix}
\subsection*{Appendix A: Mass-Polarization Term in Classical Molecular Hamiltonian}

We start from the laboratory-frame kinetic energy

\begin{equation}
T
=
\frac{1}{2}\sum_j M_j \dot{\mathbf R}_j^2
+
\frac{1}{2}\sum_i m_e \dot{\mathbf r}_i^2 ,
\end{equation}
where \(M_j\) is the mass of the \(j\)-th nucleus, \(\mathbf R_j\) is its Cartesian coordinate, \(m_e\) is the electron mass, and \(\mathbf r_i\) is the coordinate of the \(i\)-th electron.

\subsubsection{Nuclear Center-of-Mass Coordinate}

Define the nuclear total mass

\begin{equation}
M_N = \sum_j M_j .
\end{equation}
The nuclear center-of-mass coordinate is
\begin{equation}
\mathbf{R}_{\mathrm{cm}}
=
\mathbf R_N^{\rm CoM}
=
\frac{1}{M_N}\sum_j M_j \mathbf R_j .
\end{equation}
We introduce coordinates shifted by the nuclear center of mass:
\begin{equation}
\mathbf R_j'
=
\mathbf R_j - \mathbf{R}_{\mathrm{cm}} ,
\end{equation}

\begin{equation}
\mathbf r_i'
=
\mathbf r_i - \mathbf{R}_{\mathrm{cm}} .
\end{equation}
Therefore,
\begin{equation}
\mathbf R_j = \mathbf{R}_{\mathrm{cm}} + \mathbf R_j',
\qquad
\mathbf r_i = \mathbf{R}_{\mathrm{cm}} + \mathbf r_i' .
\end{equation}
The shifted nuclear coordinates satisfy the constraint
\begin{equation}
\sum_j M_j \mathbf R_j' = 0 ,
\end{equation}
and therefore
\begin{equation}
\sum_j M_j \dot{\mathbf R}_j' = 0 .
\end{equation}
This is important: the \(\mathbf R_j'\) are not all independent.

\subsubsection{Kinetic Energy in Velocity Form}

Substitute
\begin{equation}
\dot{\mathbf R}_j
=
\dot{\mathbf{R}}_{\mathrm{cm}}
+
\dot{\mathbf R}_j',
\qquad
\dot{\mathbf r}_i
=
\dot{\mathbf{R}}_{\mathrm{cm}}
+
\dot{\mathbf r}_i'
\end{equation}
into the kinetic energy.
For the nuclear part,
\begin{align}
T_N
&=
\frac{1}{2}\sum_j M_j
\left(
\dot{\mathbf{R}}_{\mathrm{cm}}
+
\dot{\mathbf R}_j'
\right)^2
\\
&=
\frac{1}{2}M_N \dot{\mathbf{R}}_{\mathrm{cm}}^2
+
\dot{\mathbf{R}}_{\mathrm{cm}}\cdot
\sum_j M_j \dot{\mathbf R}_j'
+
\frac{1}{2}\sum_j M_j \dot{\mathbf R}_j'^2 .
\end{align}
Using
\begin{equation}
\sum_j M_j \dot{\mathbf R}_j'=0 ,
\end{equation}
we get
\begin{equation}
T_N
=
\frac{1}{2}M_N \dot{\mathbf{R}}_{\mathrm{cm}}^2
+
\frac{1}{2}\sum_j M_j \dot{\mathbf R}_j'^2 .
\end{equation}
For the electronic part,
\begin{align}
T_e
&=
\frac{1}{2}\sum_i m_e
\left(
\dot{\mathbf{R}}_{\mathrm{cm}}
+
\dot{\mathbf r}_i'
\right)^2
\\
&=
\frac{1}{2}N_{\mathrm{elec}} m_e \dot{\mathbf{R}}_{\mathrm{cm}}^2
+
m_e \dot{\mathbf{R}}_{\mathrm{cm}}\cdot \sum_i \dot{\mathbf r}_i'
+
\frac{1}{2}\sum_i m_e \dot{\mathbf r}_i'^2 ,
\end{align}
where \(N_{\mathrm{elec}}\) is the number of electrons.
Hence the total kinetic energy is
\begin{equation}
T
=
\frac{1}{2}
\left(
M_N + N_{\mathrm{elec}} m_e
\right)
\dot{\mathbf{R}}_{\mathrm{cm}}^2
+
m_e \dot{\mathbf{R}}_{\mathrm{cm}}\cdot
\sum_i \dot{\mathbf r}_i'
+
\frac{1}{2}\sum_j M_j \dot{\mathbf R}_j'^2
+
\frac{1}{2}\sum_i m_e \dot{\mathbf r}_i'^2 .
\end{equation}

\subsubsection{Canonical Momenta}

The momentum conjugate to the nuclear center-of-mass coordinate \(\mathbf{R}_{\mathrm{cm}}\) is

\begin{align}
\mathbf P_{\mathrm{cm}}
&=
\frac{\partial \mathcal{L}}{\partial \dot{\mathbf{R}}_{\mathrm{cm}}}
\\
&=
\left(
M_N + N_{\mathrm{elec}} m_e
\right)\dot{\mathbf{R}}_{\mathrm{cm}}
+
m_e \sum_i \dot{\mathbf r}_i' .
\end{align}
Since
\begin{equation}
\mathbf p_i'
=
\frac{\partial \mathcal{L}}{\partial \dot{\mathbf r}_i'}
=
m_e
\left(
\dot{\mathbf{R}}_{\mathrm{cm}}
+
\dot{\mathbf r}_i'
\right),
\end{equation}
we can write
\begin{equation}
\sum_i \mathbf p_i'
=
N_{\mathrm{elec}} m_e \dot{\mathbf{R}}_{\mathrm{cm}}
+
m_e \sum_i \dot{\mathbf r}_i' .
\end{equation}
Therefore,
\begin{equation}
\boxed{
\mathbf P_{\mathrm{cm}}
=
M_N \dot{\mathbf{R}}_{\mathrm{cm}}
+
\sum_i \mathbf p_i'
}
\end{equation}
or equivalently,
\begin{equation}
\boxed{
M_N \dot{\mathbf{R}}_{\mathrm{cm}}
=
\mathbf P_{\mathrm{cm}}
-
\sum_i \mathbf p_i'
}
\end{equation}
and hence
\begin{equation}
\boxed{
\dot{\mathbf{R}}_{\mathrm{cm}}
=
\frac{1}{M_N}
\left(
\mathbf P_{\mathrm{cm}}
-
\sum_i \mathbf p_i'
\right).
}
\end{equation}
The electron momentum conjugate to \(\mathbf r_i'\) is
\begin{equation}
\boxed{
\mathbf p_i'
=
m_e
\left(
\dot{\mathbf{R}}_{\mathrm{cm}}
+
\dot{\mathbf r}_i'
\right).
}
\end{equation}
Thus \(\mathbf p_i'\) is the lab-frame electron momentum, not simply
\(m_e \dot{\mathbf r}_i'\).
Equivalently,
\begin{equation}
m_e \dot{\mathbf r}_i'
=
\mathbf p_i'
-
m_e \dot{\mathbf{R}}_{\mathrm{cm}}.
\end{equation}
The nuclear shifted momenta are
\begin{equation}
\boxed{
\mathbf P_j'
=
M_j \dot{\mathbf R}_j' .
}
\end{equation}
Because the shifted nuclear coordinates satisfy
\begin{equation}
\sum_j M_j \mathbf R_j' = 0 ,
\end{equation}
their conjugate momenta satisfy
\begin{equation}
\boxed{
\sum_j \mathbf P_j' = 0 .
}
\end{equation}
In terms of the laboratory-frame nuclear momenta
\begin{equation}
\mathbf P_j = M_j \dot{\mathbf R}_j ,
\end{equation}
one has
\begin{equation}
\boxed{
\mathbf P_j'
=
\mathbf P_j
-
\frac{M_j}{M_N}\mathbf P_N ,
}
\end{equation}
where
\begin{equation}
\mathbf P_N
=
\sum_j \mathbf P_j
=
M_N \dot{\mathbf{R}}_{\mathrm{cm}}.
\end{equation}
Using
\begin{equation}
\mathbf P_N
=
\mathbf P_{\mathrm{cm}}
-
\sum_i \mathbf p_i',
\end{equation}
we may also write
\begin{equation}
\boxed{
\mathbf P_j'
=
\mathbf P_j
-
\frac{M_j}{M_N}
\left(
\mathbf P_{\mathrm{cm}}
-
\sum_i \mathbf p_i'
\right).
}
\end{equation}

\subsubsection{Kinetic Energy in Momentum Form}
The original kinetic energy can be written as
\begin{equation}
T
=
T_{\rm nuc,int}
+
T_{\rm nuc,CoM}
+
T_e .
\end{equation}
The nuclear internal kinetic energy is
\begin{equation}
T_{\rm nuc,int}
=
\sum_j \frac{\mathbf P_j'^2}{2M_j}.
\end{equation}
The electron kinetic energy is
\begin{equation}
T_e
=
\sum_i \frac{\mathbf p_i'^2}{2m_e}.
\end{equation}
The nuclear center-of-mass kinetic energy is
\begin{equation}
T_{\rm nuc,CoM}
=
\frac{1}{2}M_N \dot{\mathbf{R}}_{\mathrm{cm}}^2 .
\end{equation}
Using
\begin{equation}
M_N \dot{\mathbf{R}}_{\mathrm{cm}}
=
\mathbf P_{\mathrm{cm}}
-
\sum_i \mathbf p_i',
\end{equation}
we get
\begin{equation}
T_{\rm nuc,CoM}
=
\frac{1}{2M_N}
\left(
\mathbf P_{\mathrm{cm}}
-
\sum_i \mathbf p_i'
\right)^2 .
\end{equation}
Therefore the kinetic energy in terms of the new momenta is
\begin{equation}
\boxed{
T
=
\sum_j \frac{\mathbf P_j'^2}{2M_j}
+
\sum_i \frac{\mathbf p_i'^2}{2m_e}
+
\frac{1}{2M_N}
\left(
\mathbf P_{\mathrm{cm}}
-
\sum_i \mathbf p_i'
\right)^2 .
}
\end{equation}
The accompanying constraints are
\begin{equation}
\boxed{
\sum_j M_j \mathbf R_j'=0,
\qquad
\sum_j \mathbf P_j'=0.
}
\end{equation}

\subsubsection{Expanded Form}
Expanding the center-of-mass recoil term gives
\begin{align}
T
&=
\sum_j \frac{\mathbf P_j'^2}{2M_j}
+
\sum_i \frac{\mathbf p_i'^2}{2m_e}
+
\frac{\mathbf P_{\mathrm{cm}}^2}{2M_N}
-
\frac{\mathbf P_{\mathrm{cm}}}{M_N}
\cdot
\sum_i \mathbf p_i'
+
\frac{1}{2M_N}
\left(
\sum_i \mathbf p_i'
\right)^2 .
\end{align}
Thus,
\begin{equation}
\boxed{
T
=
\sum_j \frac{\mathbf P_j'^2}{2M_j}
+
\sum_i \frac{\mathbf p_i'^2}{2m_e}
+
\frac{\mathbf P_{\mathrm{cm}}^2}{2M_N}
-
\frac{\mathbf P_{\mathrm{cm}}}{M_N}
\cdot
\sum_i \mathbf p_i'
+
\frac{1}{2M_N}
\left(
\sum_i \mathbf p_i'
\right)^2 .
}
\end{equation}

\subsubsection{Hamiltonian}
If the potential energy is translationally invariant, then
\begin{equation}
V
=
V\left(
\{\mathbf R_j'\},
\{\mathbf r_i'\}
\right).
\end{equation}
The Hamiltonian is therefore
\begin{equation}
\boxed{
H
=
\sum_j \frac{\mathbf P_j'^2}{2M_j}
+
\sum_i \frac{\mathbf p_i'^2}{2m_e}
+
\frac{1}{2M_N}
\left(
\mathbf P_{\mathrm{cm}}
-
\sum_i \mathbf p_i'
\right)^2
+
V\left(
\{\mathbf R_j'\},
\{\mathbf r_i'\}
\right).
}
\end{equation}

\subsubsection{Molecular Rest Frame}
If the whole molecule is taken to be in its rest frame, then
\begin{equation}
\mathbf P_{\mathrm{cm}} = 0 .
\end{equation}
The kinetic energy becomes
\begin{equation}
\boxed{
T_{\rm rest}
=
\sum_j \frac{\mathbf P_j'^2}{2M_j}
+
\sum_i \frac{\mathbf p_i'^2}{2m_e}
+
\frac{1}{2M_N}
\left(
\sum_i \mathbf p_i'
\right)^2 .
}
\end{equation}
The last term,
\begin{equation}
\frac{1}{2M_N}
\left(
\sum_i \mathbf p_i'
\right)^2 ,
\end{equation}
is the nuclear center-of-mass recoil term, also called the mass-polarization term.

\subsubsection{Important Interpretation}

The coordinate transformation used here shifts both nuclei and electrons by the nuclear center of mass, not by the total molecular center of mass.
Therefore,
\begin{equation}
\mathbf p_i'
=
m_e
\left(
\dot{\mathbf{R}}_{\mathrm{cm}}
+
\dot{\mathbf r}_i'
\right)
\end{equation}
is the canonical momentum conjugate to \(\mathbf r_i'\). It is equal to the electron lab-frame mechanical momentum.
However,
\begin{equation}
m_e \dot{\mathbf r}_i'
\end{equation}
is only the electron momentum relative to the moving nuclear center of mass. It is not the canonical momentum conjugate to \(\mathbf r_i'\) unless the nuclear center-of-mass motion has been fixed or eliminated.
The final kinetic energy in nuclear-CoM coordinates is therefore
\begin{equation}
\boxed{
T
=
\sum_j \frac{\mathbf P_j'^2}{2M_j}
+
\sum_i \frac{\mathbf p_i'^2}{2m_e}
+
\frac{1}{2M_N}
\left(
\mathbf P_{\mathrm{cm}}
-
\sum_i \mathbf p_i'
\right)^2 ,
}
\end{equation}
with
\begin{equation}
\boxed{
\sum_j M_j \mathbf R_j'=0,
\qquad
\sum_j \mathbf P_j'=0.
}
\end{equation}



% ------------------------------------------------------------------------------
% Reference and Cited Works
% ------------------------------------------------------------------------------

\subsection*{Appendix B: Term-by-term quantization of the rovibrational kernel}

Section~4 quantizes the classical Hamiltonian (3.22). All of the nontrivial operator-ordering work lives in one place --- the \textbf{rovibrational kernel}
\begin{equation}
T_{\mathrm{rovib}}=\tfrac12\sum_{\alpha\beta}A_\alpha\,\mu_{\alpha\beta}(Q)\,A_\beta,\qquad
A_\alpha\equiv J_\alpha-L_{\mathrm{elec},\alpha}-\pi_\alpha . \tag{B.1}
\end{equation}
This appendix expands the product $\hat A_\alpha\mu_{\alpha\beta}\hat A_\beta$ into its nine pieces and quantizes each one, writing down \emph{every} algebraic step. For each term we ask two questions: (i) what is the correct \emph{Hermitian} ordering of the operators, and (ii) what is its explicit \emph{differential-operator} form once $\hat\pi_\alpha$ is written out. Repeated Cartesian indices $\alpha,\beta,\gamma\in\{1,2,3\}$ and repeated mode indices $a,b,c$ are summed.

\subsubsection*{B.1\quad Two rules and a toolkit}

\textbf{Rule 1 (conjugate of a product).} For any two operators, $(\hat X\hat Y)^{\dagger}=\hat Y^{\dagger}\hat X^{\dagger}$ --- the order \emph{reverses}. A real function $f(Q)$ and the basic momenta are each individually Hermitian,
\begin{equation*}
f^{\dagger}=f,\qquad \hat P_b^{\dagger}=\hat P_b,\qquad \hat J_\alpha^{\dagger}=\hat J_\alpha,\qquad \hat L_\alpha^{\dagger}=\hat L_\alpha ,
\end{equation*}
and an operator $\hat O$ is \emph{Hermitian} (a legitimate piece of a kinetic-energy operator) when $\hat O^{\dagger}=\hat O$.

\textbf{Rule 2 (sliding operators past each other).} $\hat X\hat Y=\hat Y\hat X+[\hat X,\hat Y]$. We use this to move a momentum to the other side of a $Q$-dependent function, picking up the commutator as a correction.

\textbf{Toolkit.} The three momenta live in three separate ``sectors'': $\hat J_\alpha$ acts on the Euler angles $\Omega$, $\hat L_\alpha\equiv\hat L_{\mathrm{elec},\alpha}$ on the electron coordinates $\bar{\mathbf r}_i$, and $\hat P_b=-i\hbar\,\partial_{Q_b}$ on the normal coordinates $Q$. \emph{Operators from different sectors commute.} Explicitly [(4.7)--(4.9)]:
\begin{align}
&[\hat J_\alpha,\hat J_\beta]=-i\hbar\,\epsilon_{\alpha\beta\gamma}\hat J_\gamma,\qquad
[\hat L_\alpha,\hat L_\beta]=+i\hbar\,\epsilon_{\alpha\beta\gamma}\hat L_\gamma,\qquad
[\hat J_\alpha,\hat L_\beta]=0, \tag{B.2}\\[2pt]
&[\hat J_\alpha,f(Q)]=0,\quad [\hat L_\alpha,f(Q)]=0,\quad [\hat J_\alpha,\hat P_b]=0,\quad [\hat L_\alpha,\hat P_b]=0, \tag{B.3}\\[2pt]
&[\hat P_b,f(Q)]=-i\hbar\,\partial_{Q_b}f
\;\;\Longrightarrow\;\;
[\hat P_b,\mu_{\alpha\beta}]=-i\hbar\,\partial_b\mu_{\alpha\beta},\quad
[\hat P_b,\sigma_{c\alpha}]=-i\hbar\,\zeta^{\alpha}_{bc}. \tag{B.4}
\end{align}
The last line uses the vibrational angular momentum of (3.17a)--(3.20a),
\begin{equation*}
\hat\pi_\alpha=\sum_b\sigma_{b\alpha}(Q)\,\hat P_b,\qquad
\sigma_{b\alpha}=\sum_a\zeta^{\alpha}_{ab}\,Q_a,\qquad
\zeta^{\alpha}_{ab}=\Bigl(\textstyle\sum_j\mathbf l_{ja}\times\mathbf l_{jb}\Bigr)_\alpha ,
\end{equation*}
which is \emph{linear} in $Q$ with \emph{constant} coefficients $\zeta^{\alpha}_{ab}$, so $\partial_{Q_b}\sigma_{c\alpha}=\zeta^{\alpha}_{bc}$. Because the cross product is antisymmetric, $\zeta^{\alpha}_{ab}=-\zeta^{\alpha}_{ba}$, and therefore
\begin{equation}
\zeta^{\alpha}_{bb}=0,\qquad \sum_b\partial_{Q_b}\sigma_{b\alpha}=\sum_b\zeta^{\alpha}_{bb}=0 . \tag{B.5}
\end{equation}
Finally, $\mu_{\alpha\beta}(Q)=\mu_{\beta\alpha}(Q)$ is a \emph{symmetric}, real function of $Q$ alone.

\textbf{$\hat\pi_\alpha$ is Hermitian.} This one fact is used again and again below. By Rule~1, then Rule~2 with (B.4), then (B.5):
\begin{equation}
\hat\pi_\alpha^{\dagger}=\Bigl(\sum_b\sigma_{b\alpha}\hat P_b\Bigr)^{\!\dagger}
=\sum_b\hat P_b\,\sigma_{b\alpha}
=\sum_b\bigl(\sigma_{b\alpha}\hat P_b+[\hat P_b,\sigma_{b\alpha}]\bigr)
=\hat\pi_\alpha-i\hbar\sum_b\zeta^{\alpha}_{bb}
=\hat\pi_\alpha . \tag{B.6}
\end{equation}
Had $\sum_b\zeta^{\alpha}_{bb}$ been nonzero, $\hat\pi_\alpha$ would have needed an explicit symmetrization; the antisymmetry of $\zeta$ spares us.

\subsubsection*{B.2\quad The nine terms}

Substitute $\hat A_\alpha=\hat J_\alpha-\hat L_\alpha-\hat\pi_\alpha$ into $\hat A_\alpha\mu_{\alpha\beta}\hat A_\beta$ and multiply out, \emph{keeping the operator order} (the factors need not commute):
\begin{align}
\hat A_\alpha\mu_{\alpha\beta}\hat A_\beta
=\;&\underbrace{\hat J_\alpha\mu_{\alpha\beta}\hat J_\beta}_{\text{(T1)}}
+\underbrace{\hat L_\alpha\mu_{\alpha\beta}\hat L_\beta}_{\text{(T2)}}
+\underbrace{\hat\pi_\alpha\mu_{\alpha\beta}\hat\pi_\beta}_{\text{(T3)}}\nonumber\\
&-\underbrace{\bigl(\hat J_\alpha\mu_{\alpha\beta}\hat L_\beta+\hat L_\alpha\mu_{\alpha\beta}\hat J_\beta\bigr)}_{\text{(T4) rotation--electron}}
-\underbrace{\bigl(\hat J_\alpha\mu_{\alpha\beta}\hat\pi_\beta+\hat\pi_\alpha\mu_{\alpha\beta}\hat J_\beta\bigr)}_{\text{(T5) rotation--vibration}}\nonumber\\
&+\underbrace{\bigl(\hat L_\alpha\mu_{\alpha\beta}\hat\pi_\beta+\hat\pi_\alpha\mu_{\alpha\beta}\hat L_\beta\bigr)}_{\text{(T6) vibration--electron}} . \tag{B.7}
\end{align}
The three diagonal terms (T1)--(T3) come with the factor $\tfrac12$ of (B.1); each cross-coupling (T4)--(T6) collects the two equal classical contributions $X\mu Y$ and $Y\mu X$, which is why it appears as a symmetric \emph{sum} of orderings. We treat them one at a time.

\subsubsection*{B.3\quad Diagonal terms}

\paragraph{(T1)\ \ $\hat J_\alpha\mu_{\alpha\beta}\hat J_\beta$.}
\emph{Is it Hermitian?} Apply Rule~1 ($\mu$ real, $\hat J$ Hermitian), then relabel the dummy indices $\alpha\leftrightarrow\beta$, then use $\mu_{\beta\alpha}=\mu_{\alpha\beta}$:
\begin{equation*}
\bigl(\hat J_\alpha\mu_{\alpha\beta}\hat J_\beta\bigr)^{\dagger}
=\hat J_\beta\,\mu_{\alpha\beta}\,\hat J_\alpha
=\hat J_\alpha\,\mu_{\beta\alpha}\,\hat J_\beta
=\hat J_\alpha\,\mu_{\alpha\beta}\,\hat J_\beta .
\end{equation*}
So (T1) is Hermitian exactly as written. \emph{Ordering / extra term?} Since $\mu$ is a function of $Q$ only, it commutes with $\hat J$ [(B.3)]: $\hat J_\alpha\mu_{\alpha\beta}\hat J_\beta=\mu_{\alpha\beta}\hat J_\alpha\hat J_\beta$. Split $\hat J_\alpha\hat J_\beta=\tfrac12\{\hat J_\alpha,\hat J_\beta\}+\tfrac12[\hat J_\alpha,\hat J_\beta]$ and use (B.2):
\begin{equation*}
\mu_{\alpha\beta}\hat J_\alpha\hat J_\beta
=\tfrac12\mu_{\alpha\beta}\{\hat J_\alpha,\hat J_\beta\}
-\tfrac{i\hbar}{2}\,\mu_{\alpha\beta}\,\epsilon_{\alpha\beta\gamma}\,\hat J_\gamma .
\end{equation*}
The commutator term dies because $\mu_{\alpha\beta}$ is symmetric while $\epsilon_{\alpha\beta\gamma}$ is antisymmetric in $\alpha\beta$, so the contraction $\mu_{\alpha\beta}\epsilon_{\alpha\beta\gamma}=0$. Hence
\begin{equation}
\hat J_\alpha\mu_{\alpha\beta}\hat J_\beta=\tfrac12\,\mu_{\alpha\beta}\{\hat J_\alpha,\hat J_\beta\}
\qquad(\text{Hermitian; no }c\text{-number}). \tag{B.8}
\end{equation}
This is the asymmetric-top kinetic operator with $Q$-dependent inertia. It contains no $\hat P$, hence no differential-operator subtlety.

\paragraph{(T2)\ \ $\hat L_\alpha\mu_{\alpha\beta}\hat L_\beta$.}
Word-for-word the same as (T1) with $\hat J\to\hat L$: $\mu$ commutes with $\hat L$ [(B.3)], and although the $\hat L$ algebra carries the opposite sign, the commutator term again vanishes against symmetric $\mu$. Thus
\begin{equation}
\hat L_\alpha\mu_{\alpha\beta}\hat L_\beta=\tfrac12\,\mu_{\alpha\beta}\{\hat L_\alpha,\hat L_\beta\}
\qquad(\text{Hermitian; no }c\text{-number}). \tag{B.9}
\end{equation}

\paragraph{(T3)\ \ $\hat\pi_\alpha\mu_{\alpha\beta}\hat\pi_\beta$.}
\emph{Is it Hermitian?} Use $\hat\pi^{\dagger}=\hat\pi$ from (B.6), then relabel $\alpha\leftrightarrow\beta$ and use $\mu$ symmetric:
\begin{equation*}
\bigl(\hat\pi_\alpha\mu_{\alpha\beta}\hat\pi_\beta\bigr)^{\dagger}
=\hat\pi_\beta\,\mu_{\alpha\beta}\,\hat\pi_\alpha
=\hat\pi_\alpha\,\mu_{\beta\alpha}\,\hat\pi_\beta
=\hat\pi_\alpha\,\mu_{\alpha\beta}\,\hat\pi_\beta .
\end{equation*}
So (T3) is Hermitian as written --- but \emph{only} because $\hat\pi$ is Hermitian, i.e.\ ultimately because $\zeta$ is antisymmetric (B.5). \emph{Differential-operator form.} Write $\hat\pi_\alpha=\sum_b\sigma_{b\alpha}\hat P_b$ and slide the left momentum $\hat P_b$ rightward past the $Q$-function $\mu_{\alpha\beta}\sigma_{c\beta}$ with Rule~2 and (B.4):
\begin{align}
\hat\pi_\alpha\mu_{\alpha\beta}\hat\pi_\beta
&=\sum_{bc}\sigma_{b\alpha}\,\hat P_b\,\bigl(\mu_{\alpha\beta}\sigma_{c\beta}\bigr)\,\hat P_c
=\sum_{bc}\sigma_{b\alpha}\Bigl[\bigl(\mu_{\alpha\beta}\sigma_{c\beta}\bigr)\hat P_b-i\hbar\,\partial_b\!\bigl(\mu_{\alpha\beta}\sigma_{c\beta}\bigr)\Bigr]\hat P_c\nonumber\\[2pt]
&=\underbrace{\sum_{bc}\bigl(\boldsymbol\sigma\boldsymbol\mu\boldsymbol\sigma^{\mathsf T}\bigr)_{bc}\,\hat P_b\hat P_c}_{\text{2nd order in }\hat P}
\;-\;i\hbar\underbrace{\sum_{bc}\sigma_{b\alpha}\,\partial_b\!\bigl(\mu_{\alpha\beta}\sigma_{c\beta}\bigr)\,\hat P_c}_{\text{1st order in }\hat P}, \tag{B.10}
\end{align}
with $(\boldsymbol\sigma\boldsymbol\mu\boldsymbol\sigma^{\mathsf T})_{bc}=\sum_{\alpha\beta}\sigma_{b\alpha}\mu_{\alpha\beta}\sigma_{c\beta}$, the Coriolis-dressed vibrational inertia. The two pieces together \emph{are} the Hermitian operator (T3); note that every surviving term still carries at least one $\hat P$, so there is \emph{no} pure constant ($c$-number).

\subsubsection*{B.4\quad Cross terms}

\paragraph{(T4)\ \ rotation--electron, $\hat J_\alpha\mu_{\alpha\beta}\hat L_\beta+\hat L_\alpha\mu_{\alpha\beta}\hat J_\beta$.}
Here all three factors mutually commute: $[\hat J,\hat L]=0$, and $\mu$ commutes with both [(B.2)--(B.3)]. A single ordering is already Hermitian,
\begin{equation*}
\bigl(\hat J_\alpha\mu_{\alpha\beta}\hat L_\beta\bigr)^{\dagger}
=\hat L_\beta\mu_{\alpha\beta}\hat J_\alpha
=\mu_{\alpha\beta}\hat L_\beta\hat J_\alpha
=\mu_{\alpha\beta}\hat J_\alpha\hat L_\beta
=\hat J_\alpha\mu_{\alpha\beta}\hat L_\beta ,
\end{equation*}
and the two orderings in the sum are equal ($\hat L_\alpha\mu_{\alpha\beta}\hat J_\beta\overset{\alpha\leftrightarrow\beta}{=}\hat L_\beta\mu_{\alpha\beta}\hat J_\alpha=\hat J_\alpha\mu_{\alpha\beta}\hat L_\beta$), so
\begin{equation}
\hat J_\alpha\mu_{\alpha\beta}\hat L_\beta+\hat L_\alpha\mu_{\alpha\beta}\hat J_\beta
=2\,\mu_{\alpha\beta}\,\hat J_\alpha\hat L_\beta
\qquad(\text{Hermitian; no }c\text{-number}). \tag{B.11}
\end{equation}
This is the rotation--electron Coriolis coupling. No $\hat P$ appears, so there is no differential subtlety.

\paragraph{(T5)\ \ rotation--vibration, $\hat J_\alpha\mu_{\alpha\beta}\hat\pi_\beta+\hat\pi_\alpha\mu_{\alpha\beta}\hat J_\beta$.}
Now the subtlety: $\hat\pi$ does \emph{not} commute with $\mu$,
\begin{equation*}
[\hat\pi_\beta,\mu_{\alpha\beta}]=\sum_b\sigma_{b\beta}[\hat P_b,\mu_{\alpha\beta}]=-i\hbar\sum_b\sigma_{b\beta}\partial_b\mu_{\alpha\beta}\neq0 ,
\end{equation*}
so the single ordering $\hat J_\alpha\mu_{\alpha\beta}\hat\pi_\beta$ is \emph{not} Hermitian by itself. Its conjugate is, however, exactly the \emph{other} ordering (use $\hat J$ commutes with $\mu,\hat\pi$, then relabel $\alpha\leftrightarrow\beta$):
\begin{equation*}
\bigl(\hat J_\alpha\mu_{\alpha\beta}\hat\pi_\beta\bigr)^{\dagger}
=\hat\pi_\beta\mu_{\alpha\beta}\hat J_\alpha
\overset{\alpha\leftrightarrow\beta}{=}\hat\pi_\alpha\mu_{\alpha\beta}\hat J_\beta .
\end{equation*}
The cross term is therefore of the form $\hat X+\hat X^{\dagger}$, which is automatically Hermitian. \emph{This is why the classical product $-2\,J\mu\pi$ must be quantized as the symmetric sum, not as one ordering.} Pulling $\hat J$ out (it commutes with both $\mu$ and $\hat\pi$),
\begin{equation}
\hat J_\alpha\mu_{\alpha\beta}\hat\pi_\beta+\hat\pi_\alpha\mu_{\alpha\beta}\hat J_\beta
=\hat J_\alpha\,\{\mu_{\alpha\beta},\hat\pi_\beta\},
\qquad \{\mu_{\alpha\beta},\hat\pi_\beta\}\equiv\mu_{\alpha\beta}\hat\pi_\beta+\hat\pi_\beta\mu_{\alpha\beta} , \tag{B.12}
\end{equation}
where $\{\mu_{\alpha\beta},\hat\pi_\beta\}$, being the anticommutator of two Hermitian operators, is manifestly Hermitian. \emph{Differential-operator form.} Expand the anticommutator with $\hat\pi_\beta=\sum_c\sigma_{c\beta}\hat P_c$, sliding $\hat P_c$ rightward in the second piece [(B.4)]:
\begin{align}
\{\mu_{\alpha\beta},\hat\pi_\beta\}
&=\sum_c\Bigl(\mu_{\alpha\beta}\sigma_{c\beta}\hat P_c+\sigma_{c\beta}\hat P_c\,\mu_{\alpha\beta}\Bigr)
=\sum_c\Bigl(\underbrace{2\,\mu_{\alpha\beta}\sigma_{c\beta}\hat P_c}_{\text{1st order in }\hat P}
\;\underbrace{-\,i\hbar\,\sigma_{c\beta}\partial_c\mu_{\alpha\beta}}_{\text{0th order in }\hat P}\Bigr). \tag{B.13}
\end{align}
Multiplied by $\hat J_\alpha$, the first piece is the rotation--vibration Coriolis coupling ($\hat J$ times $\hat P$); the second is the imaginary, $\hat J$-linear term required to make the sum Hermitian (it is not a stand-alone $c$-number: it still carries $\hat J_\alpha$). Indeed one verifies Hermiticity using (B.5): $\bigl(2\mu_{\alpha\beta}\sigma_{c\beta}\hat P_c\bigr)^{\dagger}=2\mu_{\alpha\beta}\sigma_{c\beta}\hat P_c-2i\hbar\,\partial_c(\mu_{\alpha\beta}\sigma_{c\beta})$ and $\sum_c\partial_c\sigma_{c\beta}=0$, so the extra $-i\hbar\sigma_{c\beta}\partial_c\mu_{\alpha\beta}$ is exactly what (B.13) supplies.

\paragraph{(T6)\ \ vibration--electron, $\hat L_\alpha\mu_{\alpha\beta}\hat\pi_\beta+\hat\pi_\alpha\mu_{\alpha\beta}\hat L_\beta$.}
Identical to (T5) with $\hat J\to\hat L$ ($\hat L$ also commutes with $\mu$ and $\hat\pi$):
\begin{equation}
\hat L_\alpha\mu_{\alpha\beta}\hat\pi_\beta+\hat\pi_\alpha\mu_{\alpha\beta}\hat L_\beta
=\hat L_\alpha\,\{\mu_{\alpha\beta},\hat\pi_\beta\}, \tag{B.14}
\end{equation}
with the same expansion (B.13). This is the vibration--electron angular coupling.

\subsubsection*{B.5\quad Assembly, and where the pseudopotential is (and is not)}

Collecting (B.8)--(B.14), the quantized kernel $\tfrac12\hat A_\alpha\mu_{\alpha\beta}\hat A_\beta$ is Hermitian \emph{exactly} as written in (4.10)/(4.14), with $\mu$ sandwiched between the two factors. Two lessons stand out.
\begin{itemize}
\item Terms made only of \emph{mutually commuting}, individually Hermitian operators --- (T1) $\hat J\mu\hat J$, (T2) $\hat L\mu\hat L$, and (T4) $\hat J\mu\hat L$ --- are Hermitian in a \emph{single} ordering; the asymmetric-top operators (T1),(T2) shed their commutator pieces against symmetric $\mu$.
\item Terms pairing $\hat\pi$ with $\hat J$ or $\hat L$ --- (T5),(T6) --- are Hermitian only as the \emph{symmetric sum} (anticommutator), because $\hat\pi$ and $\mu$ do not commute. Term (T3) $\hat\pi\mu\hat\pi$ is the fortunate middle case: with $\mu$ between two \emph{equal} Hermitian factors it is Hermitian automatically.
\end{itemize}
\textbf{No $c$-number from the kernel.} In the explicit differential forms (B.10) and (B.13), every surviving term carries at least one momentum ($\hat P$, $\hat J$, or $\hat L$): the ordering of the kernel produces only second- and first-order operators, never a pure constant. \emph{Therefore the kernel ordering does not generate the Watson pseudopotential.} The pseudopotential $\hat U=-\tfrac{\hbar^2}{8}\sum_\alpha\mu_{\alpha\alpha}$ of (4.11) is a separate, $\mathcal O(\hbar^2)$ pure $c$-number that originates in the \emph{curved integration measure} $\sqrt{|g|}$ --- the Podolsky / Laplace--Beltrami construction --- and is extracted in \S\,4.4 by the rescaling $\psi_W=\gamma^{1/4}\psi$. The division of labour is clean: \emph{ordering of the kernel} $\rightarrow$ the Hermitian operators of this appendix; \emph{the measure} $\rightarrow$ the pseudopotential of \S\,4.4.


\newpage

\subsection*{Appendix C: Term-by-term expansion of the Podolsky operator, with every derivative}

Section~4 quantized the classical Hamiltonian (3.22) by applying the Podolsky operator (4.5) to the ro-vibrational metric (4.15) and rescaling to the flat measure. This appendix supplies what the main text compressed: the expansion of the quantized kernel of (4.41),
\begin{equation}
\hat T_{\mathrm{kernel}}=\tfrac12\sum_{\alpha\beta}\hat A_\alpha\,\mu_{\alpha\beta}(Q)\,\hat A_\beta,\qquad
\hat A_\alpha\equiv\hat J_\alpha-\hat L_\alpha-\hat\pi_\alpha,\qquad
\hat L_\alpha\equiv\hat L_{\mathrm{elec},\alpha}, \tag{C.1}
\end{equation}
into its nine products, with \emph{every} coordinate derivative written out and the Hermiticity of each block verified from first principles. All arguments use only: the adjoint of a product reverses the order, $(\hat X\hat Y)^{\dagger}=\hat Y^{\dagger}\hat X^{\dagger}$; integration by parts on the measure $d\mu_W$; and the product rule. No commutator algebra of body-fixed angular momenta appears anywhere (cf.\ the Remark in \S\,4.11). Repeated Cartesian indices $\alpha,\beta\in\{1,2,3\}$ and mode indices $b,c\in\{1,\dots,f\}$ are summed.

\subsubsection*{C.1\quad Toolkit}

The three momentum-like operators live in three separate coordinate sectors:
\begin{equation}
\hat J_\alpha=-i\hbar\,(\mathsf E^{-1})_{s\alpha}\partial_{\Omega_s}\ \ \text{[eq.~(4.24), Euler angles]},\quad
\hat P_b=-i\hbar\,\partial_{Q_b}\ \ \text{[modes]},\quad
\hat L_\alpha=-i\hbar\,\epsilon_{\alpha\beta\gamma}\textstyle\sum_i\bar r_{i\beta}\partial_{\bar r_{i\gamma}}\ \ \text{[electrons]}. \tag{C.2}
\end{equation}
Because they differentiate disjoint sets of variables, \emph{operators from different sectors commute}, and each sector's functions are constants for the other sectors' derivatives. Each is Hermitian on its factor of $d\mu_W$: $\hat J_\alpha$ on $\sin\theta\,d\Omega$ by the divergence identity (4.26) [via (4.27)], $\hat L_\alpha$ on the flat electron measure by (4.18), and $\hat P_b$ on the flat $\prod_bdQ_b$ trivially. The only non-commutativity in the problem is the product rule within the $Q$-sector,
\begin{equation}
\hat P_b\,f(Q)=f(Q)\,\hat P_b-i\hbar\,(\partial_{Q_b}f)
\qquad\Longrightarrow\qquad
\hat P_b\,\mu_{\alpha\beta}=\mu_{\alpha\beta}\hat P_b-i\hbar\,\partial_b\mu_{\alpha\beta},\quad
\hat P_b\,\sigma_{\alpha c}=\sigma_{\alpha c}\hat P_b-i\hbar\,\zeta^{\alpha}_{bc}, \tag{C.3}
\end{equation}
where the last relation uses $\sigma_{\alpha b}=\sum_a\zeta^{\alpha}_{ab}Q_a$ [(3.17a)], linear in $Q$ with \emph{constant} $\zeta^{\alpha}_{ab}=\bigl(\sum_j\mathbf l_{ja}\times\mathbf l_{jb}\bigr)_\alpha$, so $\partial_{Q_b}\sigma_{\alpha c}=\zeta^{\alpha}_{bc}$. The antisymmetry of the cross product gives the sum rule that recurs throughout:
\begin{equation}
\zeta^{\alpha}_{ab}=-\zeta^{\alpha}_{ba}\qquad\Longrightarrow\qquad
\zeta^{\alpha}_{bb}=0,\qquad \sum_b\partial_{Q_b}\sigma_{\alpha b}=\sum_b\zeta^{\alpha}_{bb}=0 . \tag{C.4}
\end{equation}
\textbf{$\hat\pi_\alpha$ is Hermitian on the flat mode measure.} With $\hat\pi_\alpha=\sum_b\sigma_{\alpha b}\hat P_b$, take the adjoint (order reverses; $\sigma$ real, $\hat P_b$ Hermitian), then slide $\hat P_b$ back to the right with (C.3), then apply (C.4):
\begin{equation}
\hat\pi_\alpha^{\dagger}
=\sum_b\hat P_b\,\sigma_{\alpha b}
=\sum_b\bigl(\sigma_{\alpha b}\hat P_b-i\hbar\,\zeta^{\alpha}_{bb}\bigr)
=\hat\pi_\alpha . \tag{C.5}
\end{equation}
Had $\sum_b\zeta^\alpha_{bb}$ been nonzero, $\hat\pi_\alpha$ would have required explicit symmetrization; the antisymmetry of $\zeta$ spares us. Finally, $\mu_{\alpha\beta}(Q)=\mu_{\beta\alpha}(Q)$ is a real symmetric function of $Q$ alone.

\subsubsection*{C.2\quad The nine products}

Substitute $\hat A_\alpha=\hat J_\alpha-\hat L_\alpha-\hat\pi_\alpha$ into (C.1) and multiply out, \emph{keeping the operator order} (the factors need not commute):
\begin{align}
\hat A_\alpha\mu_{\alpha\beta}\hat A_\beta
=\;&\underbrace{\hat J_\alpha\mu_{\alpha\beta}\hat J_\beta}_{\text{(T1) rotation}}
+\underbrace{\hat L_\alpha\mu_{\alpha\beta}\hat L_\beta}_{\text{(T2) electronic}}
+\underbrace{\hat\pi_\alpha\mu_{\alpha\beta}\hat\pi_\beta}_{\text{(T3) vibrational}}
-\underbrace{\bigl(\hat J_\alpha\mu_{\alpha\beta}\hat L_\beta+\hat L_\alpha\mu_{\alpha\beta}\hat J_\beta\bigr)}_{\text{(T4) rotation--electron}}\nonumber\\
&-\underbrace{\bigl(\hat J_\alpha\mu_{\alpha\beta}\hat\pi_\beta+\hat\pi_\alpha\mu_{\alpha\beta}\hat J_\beta\bigr)}_{\text{(T5) rotation--vibration}}
+\underbrace{\bigl(\hat L_\alpha\mu_{\alpha\beta}\hat\pi_\beta+\hat\pi_\alpha\mu_{\alpha\beta}\hat L_\beta\bigr)}_{\text{(T6) vibration--electron}} . \tag{C.6}
\end{align}
Each cross term collects the two equal classical contributions $X\mu Y$ and $Y\mu X$, which is why it appears as a symmetric sum of orderings. We now treat the six blocks one at a time: for each we verify Hermiticity and write the explicit differential form.

\subsubsection*{C.3\quad Diagonal terms}

\paragraph{(T1)\ \ $\hat J_\alpha\mu_{\alpha\beta}\hat J_\beta$.}
\emph{Hermiticity.} Take the adjoint ($\mu$ real, $\hat J_\alpha^\dagger=\hat J_\alpha$ on $\sin\theta\,d\Omega$; order reverses), then relabel the summed dummy indices $\alpha\leftrightarrow\beta$, then use $\mu_{\beta\alpha}=\mu_{\alpha\beta}$:
\begin{equation}
\bigl(\hat J_\alpha\mu_{\alpha\beta}\hat J_\beta\bigr)^{\dagger}
=\hat J_\beta\,\mu_{\alpha\beta}\,\hat J_\alpha
=\hat J_\alpha\,\mu_{\beta\alpha}\,\hat J_\beta
=\hat J_\alpha\,\mu_{\alpha\beta}\,\hat J_\beta . \tag{C.7}
\end{equation}
Hermitian exactly as written --- no symmetrization, no commutator evaluated. \emph{Differential form.} $\mu(Q)$ commutes with the Euler derivatives, and inserting (C.2),
\begin{equation}
\hat J_\alpha\mu_{\alpha\beta}\hat J_\beta
=-\hbar^{2}\mu_{\alpha\beta}\,(\mathsf E^{-1})_{s\alpha}\Bigl[(\mathsf E^{-1})_{t\beta}\,\partial_{\Omega_s}\partial_{\Omega_t}
+\bigl(\partial_{\Omega_s}(\mathsf E^{-1})_{t\beta}\bigr)\,\partial_{\Omega_t}\Bigr], \tag{C.8}
\end{equation}
with the trigonometric entries of $\mathsf E^{-1}$ read off (4.12). Every term carries at least one Euler derivative: the rotational block generates \emph{no} pure multiplicative term, in agreement with the Laplace--Beltrami computation (4.29).

\paragraph{(T2)\ \ $\hat L_\alpha\mu_{\alpha\beta}\hat L_\beta$.}
Word-for-word the same as (T1) with $\hat J\to\hat L$: the adjoint chain (C.7) uses only $\hat L_\alpha^\dagger=\hat L_\alpha$ [(4.18)] and the symmetry of $\mu$. Hermitian as written; all terms carry electron derivatives; no multiplicative residue.

\paragraph{(T3)\ \ $\hat\pi_\alpha\mu_{\alpha\beta}\hat\pi_\beta$.}
\emph{Hermiticity.} Use $\hat\pi^{\dagger}=\hat\pi$ from (C.5), then relabel $\alpha\leftrightarrow\beta$ and use $\mu$ symmetric --- the same three-step chain as (C.7). Hermitian as written, but \emph{only} because $\hat\pi$ is Hermitian, i.e.\ ultimately because of the antisymmetry (C.4). \emph{Differential form.} Write $\hat\pi_\alpha=\sigma_{\alpha b}\hat P_b$ and slide the left momentum past the $Q$-function $\mu_{\alpha\beta}\sigma_{\beta c}$ with (C.3):
\begin{align}
\hat\pi_\alpha\mu_{\alpha\beta}\hat\pi_\beta
&=\sigma_{\alpha b}\,\hat P_b\,\bigl(\mu_{\alpha\beta}\sigma_{\beta c}\bigr)\,\hat P_c
=\sigma_{\alpha b}\Bigl[\bigl(\mu_{\alpha\beta}\sigma_{\beta c}\bigr)\hat P_b-i\hbar\,\partial_b\!\bigl(\mu_{\alpha\beta}\sigma_{\beta c}\bigr)\Bigr]\hat P_c\nonumber\\[2pt]
&=\underbrace{G_{bc}\,\hat P_b\hat P_c}_{\text{2nd order}}
\;-\;i\hbar\,\underbrace{\sigma_{\alpha b}\Bigl[(\partial_b\mu_{\alpha\beta})\sigma_{\beta c}+\mu_{\alpha\beta}\,\zeta^{\beta}_{bc}\Bigr]\hat P_c}_{\text{1st order}},
\qquad G_{bc}\equiv(\boldsymbol\sigma^{\mathsf T}\boldsymbol\mu\boldsymbol\sigma)_{bc}, \tag{C.9}
\end{align}
the Coriolis-dressed vibrational operator. Every surviving term carries at least one $\hat P$: again no pure multiplicative term. \emph{Equivalence with the ordered form.} The same sum rule shows that (T3) equals the fully left-ordered operator: moving the left $\hat P_b$ past its own coefficient $\sigma_{\alpha b}$,
\begin{equation}
\hat P_b\,G_{bc}\,\hat P_c
-\hat\pi_\alpha\mu_{\alpha\beta}\hat\pi_\beta
=\bigl[\hat P_b,\sigma_{\alpha b}\bigr]\,\mu_{\alpha\beta}\sigma_{\beta c}\hat P_c
=-i\hbar\Bigl(\textstyle\sum_b\zeta^{\alpha}_{bb}\Bigr)\mu_{\alpha\beta}\sigma_{\beta c}\hat P_c=0 , \tag{C.10}
\end{equation}
so $\tfrac12\hat\pi_\alpha\mu_{\alpha\beta}\hat\pi_\beta=\tfrac12\hat P_b\,G_{bc}\,\hat P_c$ \emph{exactly}. This identity is what lets the sandwich-ordered kernel of (4.41) match the Laplace--Beltrami vibrational block (4.31), whose dressed part is precisely $\tfrac12\hat P_bG_{bc}\hat P_c$ after the rescaling (see \S\,C.5).

\subsubsection*{C.4\quad Cross terms}

\paragraph{(T4)\ \ $\hat J_\alpha\mu_{\alpha\beta}\hat L_\beta+\hat L_\alpha\mu_{\alpha\beta}\hat J_\beta$.}
All three factors act on disjoint variables and mutually commute. A single ordering is already Hermitian,
\begin{equation}
\bigl(\hat J_\alpha\mu_{\alpha\beta}\hat L_\beta\bigr)^{\dagger}
=\hat L_\beta\,\mu_{\alpha\beta}\,\hat J_\alpha
=\mu_{\alpha\beta}\,\hat J_\alpha\hat L_\beta
=\hat J_\alpha\mu_{\alpha\beta}\hat L_\beta , \tag{C.11}
\end{equation}
and the two orderings in the sum are equal (relabel $\alpha\leftrightarrow\beta$ and commute), so
\begin{equation}
\hat J_\alpha\mu_{\alpha\beta}\hat L_\beta+\hat L_\alpha\mu_{\alpha\beta}\hat J_\beta
=2\,\mu_{\alpha\beta}\,\hat J_\alpha\hat L_\beta
\qquad(\text{Hermitian; no residue}). \tag{C.12}
\end{equation}
This is the rotation--electron Coriolis coupling; no $\hat P$ appears, so there is no differential subtlety.

\paragraph{(T5)\ \ $\hat J_\alpha\mu_{\alpha\beta}\hat\pi_\beta+\hat\pi_\alpha\mu_{\alpha\beta}\hat J_\beta$.}
Now the subtlety: $\hat\pi$ does \emph{not} commute with $\mu$, since by (C.3) $\hat\pi_\beta\mu_{\alpha\beta}=\mu_{\alpha\beta}\hat\pi_\beta-i\hbar\,\sigma_{\beta b}\partial_b\mu_{\alpha\beta}\neq\mu_{\alpha\beta}\hat\pi_\beta$. A single ordering is therefore \emph{not} Hermitian by itself; its adjoint is exactly the \emph{other} ordering ($\hat J$ commutes with $\mu$ and $\hat\pi$; relabel $\alpha\leftrightarrow\beta$):
\begin{equation}
\bigl(\hat J_\alpha\mu_{\alpha\beta}\hat\pi_\beta\bigr)^{\dagger}
=\hat\pi_\beta\,\mu_{\alpha\beta}\,\hat J_\alpha
=\hat\pi_\alpha\,\mu_{\alpha\beta}\,\hat J_\beta . \tag{C.13}
\end{equation}
The cross term is of the form $\hat X+\hat X^{\dagger}$, automatically Hermitian --- \emph{this is why the classical product $-2J\mu\pi$ must be quantized as the symmetric sum}, and it is precisely the symmetric sum that the Laplace--Beltrami operator produced by itself in (4.30)/(4.35). Pulling $\hat J$ out front and expanding the anticommutator with (C.3),
\begin{equation}
\hat J_\alpha\mu_{\alpha\beta}\hat\pi_\beta+\hat\pi_\alpha\mu_{\alpha\beta}\hat J_\beta
=\hat J_\alpha\bigl(\mu_{\alpha\beta}\hat\pi_\beta+\hat\pi_\beta\mu_{\alpha\beta}\bigr)
=\hat J_\alpha\Bigl(\underbrace{2\,\mu_{\alpha\beta}\sigma_{\beta c}\hat P_c}_{\text{1st order in }\hat P}
\;\underbrace{-\,i\hbar\,\sigma_{\beta c}\,\partial_c\mu_{\alpha\beta}}_{\text{0th order in }\hat P}\Bigr). \tag{C.14}
\end{equation}
The first piece is the rotation--vibration Coriolis coupling ($\hat J$ times $\hat P$); the second is the imaginary, $\hat J$-linear counterterm that makes the sum Hermitian. It is \emph{not} a stand-alone multiplicative residue --- it still carries $\hat J_\alpha$ --- consistent with the Laplace--Beltrami result (4.35) that the Coriolis block contributes nothing to the pseudopotential. As a cross-check of (C.14) against the pre-rescaling form (4.30): expanding the second group of (4.30) by the product rule,
\begin{equation}
\tfrac{1}{\sqrt\gamma}\,\hat P_b\,\sqrt\gamma\;\sigma_{\alpha b}\mu_{\alpha\beta}\,\hat J_\beta
=\Bigl[\sigma_{\alpha b}\mu_{\alpha\beta}\hat P_b
-i\hbar\,\partial_b\!\bigl(\sigma_{\alpha b}\mu_{\alpha\beta}\bigr)
-\tfrac{i\hbar}{2}\,(\partial_b\ln\gamma)\,\sigma_{\alpha b}\mu_{\alpha\beta}\Bigr]\hat J_\beta , \tag{C.15}
\end{equation}
and conjugating by $\gamma^{1/4}$ as in \S\,4.9 cancels the $\partial_b\ln\gamma$ terms between the two groups, reproducing exactly $-\tfrac12\hat J_\alpha\{\mu_{\alpha\beta},\hat\pi_\beta\}$ of (4.35), i.e.\ $-\tfrac12\times$(C.14).

\paragraph{(T6)\ \ $\hat L_\alpha\mu_{\alpha\beta}\hat\pi_\beta+\hat\pi_\alpha\mu_{\alpha\beta}\hat L_\beta$.}
Identical to (T5) with $\hat J\to\hat L$ ($\hat L$ also commutes with $\mu$ and $\hat\pi$):
\begin{equation}
\hat L_\alpha\mu_{\alpha\beta}\hat\pi_\beta+\hat\pi_\alpha\mu_{\alpha\beta}\hat L_\beta
=\hat L_\alpha\Bigl(2\,\mu_{\alpha\beta}\sigma_{\beta c}\hat P_c-i\hbar\,\sigma_{\beta c}\,\partial_c\mu_{\alpha\beta}\Bigr), \tag{C.16}
\end{equation}
the vibration--electron angular coupling.

\subsubsection*{C.5\quad The measure residue $U_{\mathrm{Cor}}$, every derivative}

Section 4.9 derived the identity-block residue $U_{\mathrm{vib}}$ [(4.34)] in full and deferred the residue of the dressed block. Here it is. Conjugate the dressed part of the vibrational Laplace--Beltrami block (4.31) by $\gamma^{1/4}$ and act on $\psi_W$; each line is one product rule [$G_{bc}=(\boldsymbol\sigma^{\mathsf T}\boldsymbol\mu\boldsymbol\sigma)_{bc}$, symmetric]:
\begin{align}
\partial_{Q_c}\!\bigl(\gamma^{-1/4}\psi_W\bigr)&=\gamma^{-1/4}\Bigl[\partial_{Q_c}\psi_W-\tfrac14(\partial_c\ln\gamma)\psi_W\Bigr],\nonumber\\
\gamma^{1/2}\,G_{bc}\,\partial_{Q_c}\!\bigl(\gamma^{-1/4}\psi_W\bigr)&=\gamma^{1/4}\,G_{bc}\Bigl[\partial_{Q_c}\psi_W-\tfrac14(\partial_c\ln\gamma)\psi_W\Bigr],\nonumber\\
\partial_{Q_b}\Bigl(\gamma^{1/4}\,G_{bc}\bigl[\cdots\bigr]\Bigr)
&=\gamma^{1/4}\Bigl[G_{bc}\,\partial_b\partial_c\psi_W+(\partial_bG_{bc})\,\partial_c\psi_W
+\tfrac14(\partial_b\ln\gamma)G_{bc}\,\partial_c\psi_W-\tfrac14 G_{bc}(\partial_c\ln\gamma)\,\partial_b\psi_W\nonumber\\
&\qquad\quad-\tfrac14\,\partial_b\bigl(G_{bc}\,\partial_c\ln\gamma\bigr)\psi_W
-\tfrac{1}{16}\,G_{bc}(\partial_b\ln\gamma)(\partial_c\ln\gamma)\,\psi_W\Bigr]. \tag{C.17}
\end{align}
The two first-order $\ln\gamma$ terms cancel by the symmetry of $G_{bc}$ (relabel $b\leftrightarrow c$), leaving $G\partial\partial+(\partial G)\partial$ --- which is exactly $-\tfrac{2}{\hbar^2}\times\tfrac12\hat P_bG_{bc}\hat P_c$ --- plus a pure multiplicative remainder. Multiplying by $-\tfrac{\hbar^2}{2}$ (the prefactors $\gamma^{1/4}\cdot\gamma^{-1/2}\cdot\gamma^{1/4}=1$ combine to unity),
\begin{equation}
\gamma^{1/4}\,\hat T_{\mathrm{dress}}\,\gamma^{-1/4}
=\tfrac12\,\hat P_b\,G_{bc}\,\hat P_c+U_{\mathrm{Cor}},\qquad
\boxed{\;U_{\mathrm{Cor}}=\frac{\hbar^{2}}{8}\Bigl[\partial_b\bigl(G_{bc}\,\partial_c\ln\gamma\bigr)+\tfrac14\,G_{bc}\,(\partial_b\ln\gamma)(\partial_c\ln\gamma)\Bigr].\;} \tag{C.18}
\end{equation}
By (C.10), $\tfrac12\hat P_bG_{bc}\hat P_c=\tfrac12\hat\pi_\alpha\mu_{\alpha\beta}\hat\pi_\beta$, so the operator part is exactly the (T3) sandwich of the kernel (4.41), and the entire ordering debt of the dressed block is the multiplicative $U_{\mathrm{Cor}}$. Together with $U_{\mathrm{vib}}$ of (4.34) [the same computation with $G_{bc}\to\delta_{bc}$], the total residue is
\begin{equation}
\hat U=U_{\mathrm{vib}}+U_{\mathrm{Cor}}
=\frac{\hbar^{2}}{8}\Bigl[\partial_b\bigl(g^{QQ}_{bc}\,\partial_c\ln\gamma\bigr)+\tfrac14\,g^{QQ}_{bc}\,(\partial_b\ln\gamma)(\partial_c\ln\gamma)\Bigr],
\qquad g^{QQ}_{bc}=\delta_{bc}+G_{bc}, \tag{C.19}
\end{equation}
every ingredient of which reduces, via Jacobi's formula (4.36) and the linearity (4.37), to the constant $\boldsymbol\zeta$'s and the linear $\boldsymbol\sigma(Q)$; the Watson sum rules then collapse (C.19) to the trace $\hat U=-\tfrac{\hbar^2}{8}\sum_\alpha\mu_{\alpha\alpha}$ of (4.38)~\cite{watson1968}.

\textbf{The division of labour is clean.} In the explicit forms (C.8), (C.9), (C.12), (C.14), (C.16), every surviving kernel term carries at least one derivative operator ($\hat P$, $\hat J$, or $\hat L$): \emph{the ordering of the kernel never generates a pure multiplicative term}. The pseudopotential (4.38) is an $\mathcal O(\hbar^2)$ $c$-number that originates entirely in the \emph{curved integration measure} $\sqrt\gamma$, extracted by the rescaling $\psi_W=\gamma^{1/4}\psi$ --- (C.18) is its dressed-block share. Kernel ordering $\rightarrow$ the Hermitian operators of this appendix; the measure $\rightarrow$ the pseudopotential.

\subsubsection*{C.6\quad Worked example: a bent triatomic}

To make ``every derivative'' concrete, take a bent triatomic $XY_2$ (e.g.\ H${}_2$O), a $C_{2v}$ asymmetric top with $N_{\mathrm{nucl}}=3$ and $f=3N_{\mathrm{nucl}}-6=3$ normal coordinates: the symmetric stretch $Q_1$ and bend $Q_2$ (both $a_1$) and the antisymmetric stretch $Q_3$ ($b_2$). Choose the body frame as the instantaneous principal-axis system, with $a,b$ in the molecular plane and $c\perp$ to it.

\paragraph{Inertia and its derivatives.}
At equilibrium the inertia tensor is diagonal, $\mathbf I_N^{\mathrm{eq}}=\operatorname{diag}(I_a^{e},I_b^{e},I_c^{e})$, and because the equilibrium mass distribution is planar the perpendicular-axis theorem gives the exact relation
\begin{equation}
I_c^{e}=I_a^{e}+I_b^{e}. \tag{C.20}
\end{equation}
The totally symmetric coordinates $Q_1,Q_2$ preserve the $C_{2v}$ symmetry, so to linear order they only \emph{scale} the diagonal moments, whereas the antisymmetric $Q_3$ tilts the principal axes and enters $\mathbf I_N$ only at second order. Hence
\begin{equation}
I_\alpha(Q)=I_\alpha^{e}+a_\alpha^{1}Q_1+a_\alpha^{2}Q_2+\mathcal O(Q^2),
\qquad a_\alpha^{b}\equiv\frac{\partial I_\alpha}{\partial Q_b}\Big|_{\mathrm{eq}}, \tag{C.21}
\end{equation}
with the inertial derivatives $a_\alpha^{b}$ computable from the mode vectors $\mathbf l_{jb}$ evaluated at equilibrium. These are exactly the derivatives that populate the rotation--vibration Jacobian $\sqrt{\gamma(Q)}$ of (4.21).

\paragraph{Coriolis coupling.}
The only nonzero constant Coriolis matrix in $C_{2v}$ couples the two $a_1$ modes to the $b_2$ mode about the $b$ axis (the axis whose rotation maps a symmetric into an antisymmetric displacement):
\begin{equation}
\boldsymbol\sigma_1=\zeta^{b}_{13}\,Q_3\,\hat{\mathbf b},\quad
\boldsymbol\sigma_2=\zeta^{b}_{23}\,Q_3\,\hat{\mathbf b},\quad
\boldsymbol\sigma_3=-\bigl(\zeta^{b}_{13}Q_1+\zeta^{b}_{23}Q_2\bigr)\hat{\mathbf b},
\qquad \zeta^{b}_{ab}=-\zeta^{b}_{ba}, \tag{C.22}
\end{equation}
so $\boldsymbol\sigma\boldsymbol\sigma^{T}=\bigl[(\zeta^{b}_{13}Q_3)^2+(\zeta^{b}_{23}Q_3)^2+(\zeta^{b}_{13}Q_1+\zeta^{b}_{23}Q_2)^2\bigr]\hat{\mathbf b}\hat{\mathbf b}^{T}$ corrects only the $bb$ element of the inertia,
\begin{equation}
\mathbf I'(Q)=\operatorname{diag}\!\bigl(I_a(Q),\,I_b(Q)-s(Q),\,I_c(Q)\bigr),\qquad
s(Q)\equiv(\zeta^{b}_{13}Q_3)^2+(\zeta^{b}_{23}Q_3)^2+(\zeta^{b}_{13}Q_1+\zeta^{b}_{23}Q_2)^2 . \tag{C.23}
\end{equation}
Note $s(Q)=\mathcal O(Q^2)$, so it does not affect the linear-order derivatives (C.21) but does enter the pseudopotential and the exact Jacobian.

\paragraph{Determinant, volume element, and its explicit derivatives.}
With $\mathbf I'$ diagonal in (C.23),
\begin{equation}
\gamma(Q)=\det\mathbf I'=I_a(Q)\,\bigl(I_b(Q)-s(Q)\bigr)\,I_c(Q),\qquad
\sqrt{g}=\sin\theta\,\sqrt{I_a\,(I_b-s)\,I_c}, \tag{C.24}
\end{equation}
and the vibrational Jacobian carries the explicit derivatives
\begin{equation}
\partial_{Q_b}\ln\gamma
=\frac{a_a^{b}}{I_a}+\frac{a_b^{b}-\partial_{Q_b}s}{I_b-s}+\frac{a_c^{b}}{I_c},
\qquad
\partial_{Q_3}s=2\bigl[(\zeta^{b}_{13})^2+(\zeta^{b}_{23})^2\bigr]Q_3,\;\;
\partial_{Q_1}s=2\zeta^{b}_{13}\bigl(\zeta^{b}_{13}Q_1+\zeta^{b}_{23}Q_2\bigr),\ \text{etc.}, \tag{C.25}
\end{equation}
which are precisely the first-order remainders exhibited under (4.31) and removed by the $\gamma^{1/4}$ rescaling of \S\,4.9.

\paragraph{Pseudopotential.}
Since $\mathbf I'$ is diagonal, $\boldsymbol\mu=\operatorname{diag}\bigl(1/I_a,\,1/(I_b-s),\,1/I_c\bigr)$ and the trace formula (4.38) is fully explicit:
\begin{equation}
\boxed{\;\hat U(Q)=-\frac{\hbar^2}{8}\Bigl[\frac{1}{I_a(Q)}+\frac{1}{I_b(Q)-s(Q)}+\frac{1}{I_c(Q)}\Bigr].\;}\tag{C.26}
\end{equation}
Expanding to the order at which it first varies with the coordinates, and using (C.20),
\begin{equation}
\hat U(Q)=-\frac{\hbar^2}{8}\Bigl(\frac{1}{I_a^{e}}+\frac{1}{I_b^{e}}+\frac{1}{I_c^{e}}\Bigr)
+\frac{\hbar^2}{8}\sum_{b=1,2}\Bigl(\frac{a_a^{b}}{(I_a^{e})^2}+\frac{a_b^{b}}{(I_b^{e})^2}+\frac{a_c^{b}}{(I_c^{e})^2}\Bigr)Q_b+\mathcal O(Q^2). \tag{C.27}
\end{equation}
The constant is an inconsequential shift of the electronic--vibrational zero; the linear term is a genuine $\hbar^2$ force on the normal coordinates --- the physical content of the Watson pseudopotential for this molecule --- produced entirely by differentiating $\sqrt{\gamma(Q)}$ through (C.25). Every step from the metric (4.15) to the surface (C.27) used only $\partial_\phi,\partial_\theta,\partial_\chi$ and $\partial_{Q_b}$.

\subsubsection*{C.7\quad Summary}

\begin{enumerate}
\item The quantized kernel (C.1) expands into the nine products (C.6). Blocks built from mutually commuting, individually Hermitian factors --- (T1) $\hat J\mu\hat J$, (T2) $\hat L\mu\hat L$, (T4) $\hat J\mu\hat L$ --- are Hermitian in a single ordering, by the adjoint-reversal argument (C.7) alone.
\item Blocks pairing $\hat\pi$ with $\hat J$ or $\hat L$ --- (T5), (T6) --- are Hermitian only as the symmetric sum $\hat X+\hat X^\dagger$, which is exactly the ordering the Laplace--Beltrami operator produces by itself [(4.30), (4.35)]. (T3) $\hat\pi\mu\hat\pi$ is the fortunate middle case, Hermitian because $\hat\pi$ is [(C.5)], and equal to the left-ordered $\tfrac12\hat P G\hat P$ by the sum rule (C.10).
\item No kernel ordering produces a pure multiplicative term. The Watson pseudopotential comes solely from the measure: $U_{\mathrm{vib}}$ [(4.34)] plus $U_{\mathrm{Cor}}$ [(C.18)] collapse to $-\tfrac{\hbar^2}{8}\sum_\alpha\mu_{\alpha\alpha}$ [(4.38)].
\item Every Hermiticity statement above used only integration by parts, the product rule, and the divergence identities (4.18), (4.26), (C.4) --- never a commutator of body-fixed angular momenta.
\end{enumerate}


\bibliographystyle{IEEEtran}
\bibliography{References}

% ------------------------------------------------------------------------------

\end{document}
